%%%%%%%%%%%%%%%%%%%%%%%%%%%%%%%%%%%%%%%%%%%%%%%%%%%%%%%%%%%%%%%%%%%%%%
% Lớp 12 | Phương trình mặt phẳng, đường thẳng, mặt cầu | Phương trình mặt phẳng
% Dạng 2H5?1-3: Viết phương trình tổng quát mặt phẳng
% Số câu: 121  (sinh tự động bởi tach_cau.py; từ nay thầy sửa trực tiếp file này)
%%%%%%%%%%%%%%%%%%%%%%%%%%%%%%%%%%%%%%%%%%%%%%%%%%%%%%%%%%%%%%%%%%%%%%

%%%%%%%%%%%%%%%% TRẮC NGHIỆM 4 PHƯƠNG ÁN %%%%%%%%%%%%%%%%

% nguồn: Phan1_Lop_12_Hinh_hoc_Chuong_5
\begin{ex}%[2H5N1-3]%[ma:c47781ad]
	Trong không gian  $Oxyz$, điểm nào dưới đây thuộc đường thẳng $d\colon\heva{&x=3+2t\\&y=1-3t\\&z=-1+t}$?
	 \choice
		{\True $M_1\left(3;1;-1\right)$} 
		{$M_2\left(2;-3;1\right)$} 
		{$M_3\left(1;3;-1\right)$} 
		{$M_{4}\left(-3;-1;1\right)$} 
	 \loigiai{Ta có $M_1\left(3;1;-1\right)$ thuộc đường thẳng $d\colon\heva{&x=3+2t\\&y=1-3t\\&z=-1+t.}$}
\end{ex}

% nguồn: Phan1_Lop_12_Hinh_hoc_Chuong_5
\begin{ex}%[2H5N1-3]%[ma:3018e67a]
	Trong không  gian $Oxyz$, phương trình nào sau đây là phương trình của mặt phẳng $Oyz$?
	 \choice
		{$y  = 0$}
		{$z =  0$}
		{\True  $x = 0$}
		{$x + 3  = 0$}
	\loigiai{Phương  trình mặt phẳng $Oyz$ là $x = 0$.
	 }
\end{ex}

% nguồn: Phan1_Lop_12_Hinh_hoc_Chuong_5
\begin{ex}%[2H5N1-3]%[ma:b881bcd5]
	Trong không  gian $Oxyz$, mặt phẳng $(Oxy)$ có phương trình là
	\choice 
		 {$x=0$}
		{$y=0$} 
		 {$x+y=0$}
		{\True  $z=0$}
	\loigiai{Mặt phẳng $(Oxy)$  có phương trình là $z=0$.}
\end{ex}

% nguồn: Phan1_Lop_12_Hinh_hoc_Chuong_5
\begin{ex}%[2H5N1-3]%[ma:8c95ec9a]
	Trong không gian với hệ trục tọa độ $Oxyz$, phương trình mặt phẳng $\left(\alpha\right)$ đi qua điểm $A(2;-1;3)$ và có vectơ pháp tuyến  $\overrightarrow{n}=(2;3;-1)$ là
	\choice 
		 {$(\alpha)\colon 2x-y+3z+2=0$}
		{\True $(\alpha)\colon  2x+3y-z+2=0$}
		 {$(\alpha)\colon 2x+3y-z-2=0$}
		 {$(\alpha)\colon 2x-y+3z-2=0$}
	\loigiai{ 
	Phương trình mặt phẳng đi qua điểm $A(2;-1;3)$ và có vectơ pháp tuyến $\overrightarrow{n}=(2;3;-1)$ là \[(\alpha)\colon 2(x-2)+3(y+1)-(z-3)=0\Leftrightarrow 2x+3y-z+2=0.\] 
	 }
\end{ex}

% nguồn: Phan1_Lop_12_Hinh_hoc_Chuong_5
\begin{ex}%[2H5N1-3]%[ma:4a09c7d5]
	Trong không gian $Oxyz$, phương trình mặt phẳng đi qua ba  điểm $A(0;3;0)$, $B(2;0;0)$, $C(0;0;-2)$ là
	 \choice
		{$2x+3y-2z=1$} 
		{$\dfrac{x}{3}+\dfrac{y}{2}+\dfrac{z}{-1}=1$} 
		{$3x+2y-2z=1$} 
		{\True $\dfrac{x}{2}+\dfrac{y}{3}+\dfrac{z}{-2}=1$} 
	\loigiai{Mặt phẳng đi qua ba điểm $A(0;3;0)$,  $B(2;0;0)$, $C(0;0;-2)$ có phương trình mặt phẳng theo đoạn chắn là
		\[\dfrac{x}{2}+\dfrac{y}{3}+\dfrac{z}{-2}=1.\] 
	 }
\end{ex}

% nguồn: Phan1_Lop_12_Hinh_hoc_Chuong_5
\begin{ex}%[2H5N1-3]%[ma:1a08366b]
	Trong không gian  $Oxyz$, cho $A(1;1;-2)$, $B(2;0;3)$, $C(-2;4;1)$. Mặt phẳng đi qua $A$ và vuông góc với đường thẳng $BC$ có phương trình là
	 \choice
		 {\True $2x-2y+z+2=0$}
		{$x+y-2x-6=0$} 
		{$x+y-2z+2=0$} 
		{$2x+2y+z-2=0$} 
	 \loigiai{
		Đường thẳng $BC$ có một  véctơ chỉ phương là $\overrightarrow{BC}=(-4;4;-2)=-2(2;-2;1)$.\\
		Mặt phẳng đi qua $A$ và vuông góc với đường thẳng $BC$ nên mặt phẳng đó có một véctơ pháp  tuyến là $\overrightarrow{n}=(2;-2;1)$.\\
		Vậy phương trình là $2(x-1)-2(y-1)+1(z+2)=0\Leftrightarrow  2x-2y+z+2=0$.
	} 
\end{ex}

% nguồn: Phan1_Lop_12_Hinh_hoc_Chuong_5
\begin{ex}%[2H5N1-3]%[ma:162136f9]
	Trong không gian với hệ tọa độ $Oxyz$, mặt phẳng $(\alpha)$ đi qua điểm $M(-1;5;2)$ đồng  thời $(\alpha)$ có cặp vectơ chỉ phương $\overrightarrow{u}=(0;1;1)$ và $\overrightarrow{v} = (-3;-5;1)$ có phương trình là
	\choice 
		 {$y + z - 7 = 0$}
		{\True  $2x - y + z + 5 = 0$}
		 {$-x + 5y + 2z + 5 = 0$}
		{$2x +  y + z - 5 = 0$}
	 \loigiai{
		Ta  có $[\overrightarrow{u},\overrightarrow{v}] = (6;-3;3)$.\\
		Chọn $\overrightarrow{n} =  (2;-1;1)$ là 1 vectơ pháp tuyến của $(\alpha)$.\\
		Phương trình mặt phẳng qua $M(-1;5;2)$ nhận $\overrightarrow{n}  = (2;-1;1)$ làm vectơ pháp tuyến là
		$$2(x+1) - (y-5) + (z-2) = 0 \Leftrightarrow 2x - y + z + 5 =  0.$$
	 }
\end{ex}

% nguồn: Phan1_Lop_12_Hinh_hoc_Chuong_5
\begin{ex}%[2H5N1-3]%[ma:0a1beae4]
	Trong không gian $Oxyz$, cho $A(7;6;7)$, $B(12;13;6)$,  mặt phẳng $(P)$ chứa $AB$ và song song với $Ox$ có phương trình là
	\choice 
		 {\True $-y-7z+55=0$}
		 {$-y-7z+53=0$}
		 {$-y-7z+56=0$}
		{$-y-7z+35=0$} 
	 \loigiai{
		Mặt phẳng $(P)$ chứa $AB$ và song song trục $Ox$ nên có cặp véctơ chỉ phương là $\overrightarrow{AB}=(5;7;-1)$  và $\overrightarrow{i}=(1;0;0)\Rightarrow \overrightarrow{n}=\left[\overrightarrow{AB},\overrightarrow{i}\right]=(0;-1;-7)$.\\
		Mặt phẳng $(P)$ đi qua điểm $A(7;6;7)$ và có một véctơ pháp tuyến $\overrightarrow{n}=(0;-1;-7)$ có phương trình là $0(x-7)-1(y-6)-7(z-7)=0\Leftrightarrow -y-7z+55=0$. 
	} 
\end{ex}

% nguồn: Phan1_Lop_12_Hinh_hoc_Chuong_5
\begin{ex}%[2H5N1-3]%[ma:c41f7142]
	Trong không gian với hệ tọa độ $Oxyz$, cho hai điểm $A(5;-4;2)$ và $B(1;2;3)$. Viết phương trình của mặt phẳng $(P)$ đi  qua $A$ và vuông góc với đường thẳng $AB$.
	\choice 
		{\True $2x - 3y - z - 20 = 0$} 
		{$3x - y + 3z  - 25 = 0$}
		{$2x  - 3y - z + 8 = 0$}
		{$3x - y  + 3z - 13 = 0$}
	 \loigiai{
		Ta có $\overrightarrow{AB} = (-4;6;2) =  -2(2;-3;-1)$.\\
		Mặt phẳng $(P)$ đi qua $A(5;-4;2)$ nhận $\overrightarrow{n} = (2;-3;-1)$  làm vectơ pháp tuyến.\\
		Do đó, phương trình mặt phẳng $(P)$: $2(x - 5) - 3(y + 4) - 1(z - 2) = 0 \Leftrightarrow 2x - 3y - z - 20 = 0$. 
	} 
\end{ex}

% nguồn: Phan1_Lop_12_Hinh_hoc_Chuong_5
\begin{ex}%[2H5N1-3]%[ma:cef373eb]
	Trong không gian với hệ tọa độ $Oxyz$, cho điểm $M(2;-1;4)$ và mặt phẳng $(P)\colon 3x - 2y + z + 1 = 0$. Phương trình của mặt phẳng đi qua $M$ và song  song với mặt phẳng $(P)$ là
	 \choice
		{$2x - 2y + 4z - 21  = 0$}
		{\True $3x - 2y +  z - 12 = 0$}
		{$2x - 2y +  4z + 21 = 0$}
		{$3x - 2y +  z + 12 = 0$}
	\loigiai{ 
		Phương trình của mặt phẳng đi qua $M$ và song song với mặt  phẳng $(P)$ là
		$$3(x - 2)  - 2(y + 1) + 1(z - 4) = 0 \Leftrightarrow 3x - 2y + z - 12 = 0.$$
	 }
\end{ex}

% nguồn: Phan1_Lop_12_Hinh_hoc_Chuong_5
\begin{ex}%[2H5N1-3]%[ma:dea19660]
	Trong không gian $Oxyz$, mặt  phẳng tọa độ $Oxy$ có phương trình là
	 \choice
		 {$y=0$}
		{\True  $z=0$}
		 {$x=0$}
		 {$x+y+z=0$}
	\loigiai{ Mặt phẳng  toạ độ $Oxy$ có phương trình là $z=0$.}
\end{ex}

% nguồn: Phan1_Lop_12_Hinh_hoc_Chuong_5
\begin{ex}%[2H5N1-3]%[ma:cbe8d8c5]
	Trong không $Oxyz$, cho điểm $A(-3;4;-2)$ và vectơ $\vec{n}=(-2;3;-4)$. Mặt  phẳng $(P)$ đi qua điểm $A$ và nhận $\vec{n}$~làm vectơ pháp tuyến có phương trình~là
	\choice 
		 {$-3x+4y+2z+26=0$}
		{$-2x+3y-4z+29=0$} 
		{$2x-3y+4z+29=0$} 
		{\True  $2x-3y+4z+26=0$}
	\loigiai{ 
		Ta có $(P) \colon \heva{&\text{ Qua } A(-3;4;-2)\\&\text{ Vectơ pháp tuyến  }\vec{n}=(-2;3;-4).}$\\
	Suy ra $(P) \colon -2(x+3)+3(y-4)-4(z+2)=0  \Leftrightarrow (P) \colon -2x+3y-4z-26=0$.}
\end{ex}

% nguồn: Phan1_Lop_12_Hinh_hoc_Chuong_5
\begin{ex}%[2H5N1-3]%[ma:ee0bf015]
	Trong không gian với hệ tọa độ $Oxyz$, phương trình nào  sau đây là phương trình tổng quát của mặt phẳng
	 \choice
		 {$x^2+y+z+2=0$}
		 {$2 x+y^2+z+1=0$}
		 {\True $2 x+y+z+3=0$}
		{$2  x+y+z^2+4=0$}
	\loigiai{Trong không gian với hệ tọa độ $Oxyz$, phương trình $2 x+y+z+3=0$ là phương trình tổng quát của  mặt phẳng.
	 }
\end{ex}

% nguồn: Phan1_Lop_12_Hinh_hoc_Chuong_5
\begin{ex}%[2H5N1-3]%[ma:ba3f45be]
	Trong không gian $Oxyz$, mặt phẳng $(P) \colon x+y+z-3=0$ đi qua điểm  nào dưới đây?
	 \choice
		 {$M(-1;-1;-1)$}
		 {\True $N(1;1;1)$}
		 {$P(-3;0;0)$}
		 {$Q(0;0;-3)$}
	\loigiai{ 
		Thay tọa  độ điểm $N$ vào phương trình mặt phẳng $(P)$ ta có $1+1+1-3=0$.\\
		Do đó $N  \in (P)$.
	 }
\end{ex}

% nguồn: Phan1_Lop_12_Hinh_hoc_Chuong_5
\begin{ex}%[2H5N1-3]%[ma:c22fd7eb]
	Trong không  gian với hệ tọa độ $Oxyz$, cho ba điểm $A(-1;1;1)$, $B(2;1;0)$, $C(1;-1;2)$. Mặt phẳng đi qua $A$ và vuông góc với đường thẳng $BC$ có phương trình là
	 \choice
		{$3x + 2z  + 1 = 0$}
		{\True $x +  2y - 2z + 1 = 0$}
		{$x +  2y - 2z = 0$}
		{$3x + 2z  - 1 = 0$}
	 \loigiai{
		Ta có $\overrightarrow{CB} = (1;2;-2)$ là một vectơ pháp tuyến của mặt phẳng $(P)$ cần  tìm.\\
		 Vậy phương trình mặt phẳng $(P)$ là
		$$1(x + 1) + 2(y - 1)  - 2(z - 1) = 0 \Leftrightarrow x + 2y - 2z + 1 = 0.$$
	 }
\end{ex}

% nguồn: Phan1_Lop_12_Hinh_hoc_Chuong_5
\begin{ex}%[2H5N1-3]%[ma:83864525]
	Trong không gian $Oxyz$, mặt phẳng  cắt ba trục toạ độ tại ba điểm $D(3; 0; 0)$, $E(0;-2; 0)$, $G(0; 0; -7)$ có phương trình chính tắc~là
	\choice 
		 {$\dfrac{x}{3}-\dfrac{y}{2}-\dfrac{z}{7}+1=0$}
		 {$\dfrac{x}{3}+\dfrac{y}{2}+\dfrac{z}{7}=1$}
		 {\True $\dfrac{x}{3}-\dfrac{y}{2}-\dfrac{z}{7}=1$}
		 {$\dfrac{x}{3}-\dfrac{y}{2}+\dfrac{z}{7}=1$}
	\loigiai{ 
		Dễ thấy rằng phương trình mặt phẳng cắt ba trục  toạ độ tại ba điểm $D(3; 0; 0)$, $E(0; -2; 0)$, $G(0; 0; -7)$ có phương trình~là $\dfrac{x}{3}+\dfrac{y}{-2}+\dfrac{z}{-7}=1
		\Leftrightarrow  \dfrac{x}{3}-\dfrac{y}{2}-\dfrac{z}{7}=1$.
	 }
\end{ex}

% nguồn: Phan1_Lop_12_Hinh_hoc_Chuong_5
\begin{ex}%[2H5N1-3]%[ma:d1f05ec0]
	Trong không gian với hệ tọa độ $Oxyz$, phương trình nào dưới đây là phương trình mặt phẳng  đi qua điểm $M(1;2;-3)$ và có một vectơ pháp tuyến $\overrightarrow{n} = (1;-2;3)$.
	 \choice
		{\True $x - 2y + 3z  + 12 = 0$}
		{$x - 2y - 3z  - 6 = 0$}
		{$x - 2y + 3z - 12  = 0$}
		{$x - 2y - 3z  + 6 = 0$}
	\loigiai{ 
		Phương trình mặt phẳng đi qua điểm $M(1;2;-3)$ và có một vectơ pháp tuyến $\overrightarrow{n} = (1;-2;3)$  là
		$$1(x - 1) - 2(y - 2) + 3(z + 3)  = 0 \Leftrightarrow x - 2y + 3z + 12 = 0.$$
	 }
\end{ex}

% nguồn: Phan1_Lop_12_Hinh_hoc_Chuong_5
\begin{ex}%[2H5N1-3]%[ma:916364e6]
	Trong không gian $Oxyz$, cho mặt phẳng $(P)\colon 2x-y+3z-9=0$. Điểm nào dưới đây thuộc mặt phẳng $(P)$? 
	 \choice
		 {$M(0;-3;2)$}
		 {$N(0;3;-2)$}
		{$E(2;-1;3)$} 
		{\True$F(0;0;-3)$} 
	\loigiai{ 
		Thay $x=0$; $y=0$; $z=-3$ vào $2x-y+3z-9=0$ ta được $0=0$  nên điểm tọa độ điểm $F(0;0;-3)$ thuộc mặt phẳng $(P)$.
	} 
\end{ex}

% nguồn: Phan1_Lop_12_Hinh_hoc_Chuong_5
\begin{ex}%[2H5N1-3]%[ma:6a9363a0]
	Phương trình nào dưới đây là phương trình  của mặt phẳng $(Oyz)$?
	\choice 
		 {$y=0$}
		{\True $x=0$} 
		{$y-z=0$} 
		{$z=0$} 
	\loigiai{ 
		Phương  trình của mặt phẳng $(O y z)$ là $x=0$.
	} 
\end{ex}

% nguồn: Phan1_Lop_12_Hinh_hoc_Chuong_5
\begin{ex}%[2H5N1-3]%[ma:7a7121b2]
	Trong  không gian $Oxyz$, phương trình mặt phẳng $(P)$ đi qua điểm $A(1;2;-3)$ và nhận vectơ $\overrightarrow{n}=(3;2;-4)$ làm một vectơ pháp tuyến là
	 \choice
		 {$x+2y-3z-7=0$}
		 {$x+2y-3z-5=0$}
		{\True $3x+2y-4z-19=0$} 
		 {$3x+2y-4z+19=0$}
	 \loigiai{
		Phương  trình mặt phẳng $(P)$ đi qua điểm $A(1;2;-3)$ và nhận vectơ $\overrightarrow{n}=(3;2;-4)$ làm một vectơ pháp tuyến là
		$$ (P)\colon 3(x - 1) + 2(y - 2)  - 4(z + 3) = 0 \Leftrightarrow 3x + 2y - 4z - 19 = 0. $$
	} 
\end{ex}

% nguồn: Phan1_Lop_12_Hinh_hoc_Chuong_5
\begin{ex}%[2H5N1-3]%[ma:c7a16a38]
	Trong không gian $Oxyz$, mặt phẳng đi qua  các điểm $A(2;0;0)$, $B(0;-3;0)$, $C(0;0;5)$ có phương trình là
	 \choice
		 {$\dfrac{x}{5}+\dfrac{y}{-3}+\dfrac{z}{2}=1$}
		 {$\dfrac{x}{-3}+\dfrac{y}{2}+\dfrac{z}{5}=1$}
		{$\dfrac{x}{2}+\dfrac{y}{-3}+\dfrac{z}{5}=0$} 
		{\True $\dfrac{x}{2}+\dfrac{y}{-3}+\dfrac{z}{5}=1$} 
	 \loigiai{
		Phương trình mặt phẳng $(ABC)$ chắn  $3$ trục tọa độ, nên có phương trình là $\dfrac{x}{2}+\dfrac{y}{-3}+\dfrac{z}{5}=1$.
	} 
\end{ex}

% nguồn: Phan1_Lop_12_Hinh_hoc_Chuong_5
\begin{ex}%[2H5N1-3]%[ma:039cadcb]
	Trong không gian $Oxyz$, phương trình mặt phẳng  đi qua điểm $M(5; 1; 2)$ và có một vectơ pháp tuyến $\overrightarrow{n}=(1;-2; 3)$ là
	\choice 
		{$5 x+y+2  z+9=0$}
		 {$x-2 y+3 z+9=0$}
		{\True  $x-2 y+3 z-9=0$}
		{$5 x+y+2  z-9=0$}
	\loigiai{ 
		Phương trình mặt phẳng đi qua điểm $M(5; 1; 2)$ và có một vectơ pháp tuyến  $\overrightarrow{n}=(1;-2; 3)$ là
		 $$1(x-5)-2(y-1)+3(z-2)=0 \Leftrightarrow x-2 y+3 z-9=0.$$
	 }
\end{ex}

% nguồn: Phan1_Lop_12_Hinh_hoc_Chuong_5
\begin{ex}%[2H5N1-3]%[ma:0b26139a]
	Trong không gian với hệ tọa độ  $Oxyz$, mặt phẳng đi qua điểm $M(a;b;c)$ và nhận $\overrightarrow{n}=(2;-3;4)$ là vectơ pháp tuyến có phương trình là
	 \choice
		{$a(x-2)+b(y+3)+c(z-4)=0$} 
		{\True $2(x-a)-3(y-b)+4(z-c)=0$} 
		{$a(x-2)+b(y-3)+c(z-4)=0$} 
		{$2(x-a)+3(y-b)+4(z-c)=0$} 
	 \loigiai{
		Mặt phẳng  đi qua điểm $M(a;b;c)$ và nhận $\overrightarrow{n}=(2;-3;4)$ là vectơ pháp tuyến có phương trình là $$2(x-a)-3(y-b)+4(z-c)=0.$$
	} 
\end{ex}

% nguồn: Phan1_Lop_12_Hinh_hoc_Chuong_5
\begin{ex}%[2H5H1-3]%[ma:7b1d1d09]
	Trong không gian với hệ tọa độ  $Oxyz$, cho mặt phẳng $(P)\colon 2x - y + 2z + 11 = 0$ và hai điểm $A(1;0;-2)$, $B(-1;1;3)$. Mặt phẳng $(Q)$ đi qua hai điểm $A$, $B$ và vuông góc với mặt phẳng $(P)$ có phương trình là
	 \choice
		{\True $3x +  14y + 4z + 5 = 0$}
		{$2x - y + 2z - 2  = 0$}
		{$2x - y  + 2z + 2 = 0$}
		{$3x + 14y + 4z - 5  = 0$}
	\loigiai{ 
		Gọi $\overrightarrow{n_{(P)}}$, $\overrightarrow{n_{(Q)}}$ lần lượt là vectơ pháp tuyến của hai mặt phẳng $(P)$ và  $(Q)$.\\
		Ta có $\overrightarrow{AB}  = (-2;1;5)$, $\overrightarrow{n_{(P)}} = (2;-1;2)$.\\
		 Vì $(Q)$ đi qua $A$, $B$ và $(Q) \perp (P)$ nên $\overrightarrow{n_{(Q)}} \perp \overrightarrow{AB}$ và $\overrightarrow{n_{(Q)}} \perp \overrightarrow{n_{(P)}}$.\\
		Chọn $\overrightarrow{n_{(Q)}} = \left[\overrightarrow{AB}, \overrightarrow{n_{(P)}}\right]  = (3;14;4)$.\\
		Phương trình của $(Q)$ là $3(x - 1) + 14(y - 0) +  4(z + 2) = 0$ hay $3x + 14y + 4z + 5 = 0$.
	} 
\end{ex}

% nguồn: Phan1_Lop_12_Hinh_hoc_Chuong_5
\begin{ex}%[2H5H1-3]%[ma:8e76044c]
	Trong không gian $Oxyz$, mặt phẳng $(P)$ đi qua điểm $A(1;0;2)$ và vuông góc với đường thẳng $d \colon  \dfrac{x}{2}=\dfrac{y-1}{-1}=\dfrac{z+2}{3}$ có phương trình là
	\choice 
		 {$2x+y-3z-8=0$}
		 {\True $2x-y+3z-8=0$}
		 {$2x+y-3z+8=0$}
		 {$2x-y+3z+8=0$}
	 \loigiai{
		Vectơ chỉ phương của đường thẳng $d$ là $\overrightarrow{u}=(2;-1;3)$.\\ 
		Vì mặt phẳng $(P)  \perp d$ nên có vectơ pháp tuyến có tọa độ là $\overrightarrow{n}=(2;-1;3)$.\\
		Phương trình mặt phẳng $(P)$ đi qua điểm $A(1;0;2)$  và vuông góc với đường thẳng $d$ là\\
		$$2(x-1)-1(y-0)+3(z-2)=0 \Leftrightarrow 2x-y+3z-8=0.$$ 
	 }
\end{ex}

% nguồn: Phan1_Lop_12_Hinh_hoc_Chuong_5
\begin{ex}%[2H5H1-3]%[ma:bacf714a]
	Trong không gian $Oxyz$, cho hai điểm $A(2;4;1)$,  $B(-1;1;3)$ và mặt phẳng \break$(P)\colon x-3y+2z-5=0$. Lập phương trình mặt phẳng $(Q)$ đi qua hai điểm $A$, $B$ và vuông góc với mặt phẳng $(P)$.
	\choice 
		{\True $2y+3z-11=0$} 
		 {$2x-3y-11=0$}
		 {$x-3y+2z-5=0$}
		{$3y+2z-11=0$} 
	\loigiai{ 
		Ta có $\overrightarrow{AB}=(-3;-3;2)$, vectơ  pháp tuyến của mặt phẳng $(P)$ là $\vec{n}_P=(1;-3;2)$.\\
		Từ giả thiết suy ra $\vec{n}=\left[\overrightarrow{AB}, \vec{n}_P\right]=(0;8;12)$ là vectơ pháp tuyến  của $(Q)$.\\
		$(Q)$ đi qua điểm  $A(2; 4; 1)$ suy ra phương trình tổng quát của $(Q)$ là
		 \[0(x-2)+8(y-4)+12(z-1)=0\Leftrightarrow 2y+3z-11=0.\]
	 }
\end{ex}

% nguồn: Phan1_Lop_12_Hinh_hoc_Chuong_5
\begin{ex}%[2H5H1-3]%[ma:0caf34a3]
	Trong không gian với hệ tọa độ $Oxyz$, cho ba điểm $A(0;2;1)$, $B(3;0;1)$, $C(1;0;0)$.  Phương trình mặt phẳng $(ABC)$ là
	 \choice
		{$2x - 3y -  4z + 2 = 0$}
		{$4x  + 6y - 8z + 2 = 0$}
		{\True $2x + 3y - 4z - 2 =  0$}
		{$2x - 3y  - 4z - 2 = 0$}
	 \loigiai{
		Ta  có $\overrightarrow{AB} = (3;-2;0)$. $\overrightarrow{AC} = (1;-2;-1)$.\\
		Vì  mặt phẳng $(ABC)$ đi qua ba điểm $A$, $B$, $C$ nên vectơ pháp tuyến\\ $\overrightarrow{n} = [\overrightarrow{AB},\overrightarrow{AC}] = (2;3;-4)$.\\
		Phương trình mặt phẳng $(ABC)\colon 2(x-0) + 3(y-2) - 4(z-1) = 0 \Leftrightarrow 2x  + 3y - 4z - 2 = 0$.
	 }
\end{ex}

% nguồn: Phan1_Lop_12_Hinh_hoc_Chuong_5
\begin{ex}%[2H5H1-3]%[ma:0b8d6eb8]
	Cho mặt phẳng $(P)$ đi qua $A(2;-1; 4)$, $B(3; 2;-1)$ và vuông góc với mặt phẳng \break$(Q)\colon x+y+2z-3=0$. Khi  đó mặt phẳng $(P)$ có phương trình là
	 \choice
		 {$11x+7y-2z-21=0$}
		{$11x+7y+2z+21=0$} 
		 {\True $11x-7y-2z-21=0$}
		{$11x-7y+2z+21=0$} 
	 \loigiai{
		Ta có  $\overrightarrow{AB}=(1; 3;-5)$; $\overrightarrow{n}_Q=(1; 1; 2)$.\\
		Mặt  phẳng $(P)$ đi qua $A, B$ và vuông góc với mặt phẳng $(Q)$
		nên  nhận $\left[\overrightarrow{AB}, \overrightarrow{n}_Q\right]=(11;-7;-2)$ làm véctơ pháp tuyến nên có phương trình $11x-7y-2z-21=0$.
	} 
\end{ex}

% nguồn: Phan1_Lop_12_Hinh_hoc_Chuong_5
\begin{ex}%[2H5H1-3]%[ma:a7c7ff26]
	Trong không gian  với hệ tọa độ $Oxyz$, viết phương trình mặt phẳng đi qua ba điểm $A(1; 1; 4)$, $B(2; 7; 9)$, $C(0; 9; 13)$.
	 \choice
		{$2x+y+z+1=0$} 
		{\True  $x-y+z-4=0$}
		 {$2x+y-z-2=0$}
		 {$7x-2y+z-9=0$}
	 \loigiai{
		Mặt phẳng $(ABC)$ có véc-tơ pháp tuyến là $\overrightarrow{n}=\left[\overrightarrow{AB},\overrightarrow{AC}\right]=(14;-14;14)=14(1;-1;1)$  và đi qua $A(1;1;4)$ có phương trình là
	\[1(x-1)-(y-1)+1(z-4)=0\Leftrightarrow x-y+z-4=0.\]} 
\end{ex}

% nguồn: Phan1_Lop_12_Hinh_hoc_Chuong_5
\begin{ex}%[2H5H1-3]%[ma:ff39a00b]
	Trong không gian với hệ tọa độ $Oxyz$, phương trình mặt phẳng $(P)$ đi qua điểm $M(2;1;-2)$  và vuông góc với hai mặt phẳng $(Q): x + y + 2z - 3 = 0$ và $(R): x - y - z + 4 = 0$ là
	 \choice
		{$2x + 3y  - 2z + 9 = 0$}
		{$2x + 3y - 2z - 9  = 0$}
		{\True $x + 3y -  2z - 9 = 0$}
		{$2x + y - 2z + 9  = 0$}
	 \loigiai{
		Mặt phẳng $(Q)$, $(R)$ có một vectơ pháp  tuyến lần lượt là $\overrightarrow{n_{(Q)}} = (1;1;2)$ và \break $\overrightarrow{n_{(R)}} = (1;-1;-1)$.\\
		Mặt phẳng $(P)$ đi qua điểm $M(2;1;-2)$ và vuông góc với hai mặt  phẳng $(Q)$ và $(R)$ nên có cặp vectơ chỉ phương là $\overrightarrow{n_{(Q)}} = (1;1;2)$ và $\overrightarrow{n_{(R)}} = (1;-1;-1)$.\\
		Do đó $(P)$ có vectơ  pháp tuyến là
		$\overrightarrow{n_{(P)}}  = [\overrightarrow{n_{(Q)}}, \overrightarrow{n_{(R)}}] = (1;3;-2)$.\\
		Mặt phẳng $(P)$ đi qua điểm $M(2;1;-2)$ và nhận $\overrightarrow{n_{(P)}} = (1;3;-2)$ làm một vectơ pháp tuyến nên có phương trình: $1(x - 2) + 3(y -  1) - 2(z + 2) = 0 \Leftrightarrow x + 3y - 2z - 9 = 0$.
	 }
\end{ex}

% nguồn: Phan1_Lop_12_Hinh_hoc_Chuong_5
\begin{ex}%[2H5H1-3]%[ma:65ffe788]
	Trong không gian  $Oxyz$, mặt phẳng đi qua điểm $M(1;-2;3)$ và song song với mặt phẳng $(Q) \colon 2x-3y+z+5=0$ có phương trình~là
	 \choice
		{$2x-3y+z-11=0$} 
		{\True $2x-3y+z+11=0$} 
		 {$x-2y+3z-11=0$}
		 {$x-2y+3z+11=0$}
	\loigiai{ 
		Gọi $(\alpha)$~là mặt phẳng  cần tìm.\\
		Theo dữ kiện ta có $(\alpha) \colon \heva{&\text{ Qua }M(1;-2;3)\\&(\alpha)  \parallel (Q)}$\\
		Suy ra $(\alpha) \colon 2x-3y+z+d=0$ thay  $M(1;-2;3)$ vào ta được.\\
		$\Rightarrow  2\cdot1-3\cdot(-2)+3+d=0 \Leftrightarrow d=11$ (nhận do $d \ne 5$)\\
		Suy ra $(\alpha) \colon  2x-3y+z+11=0$.
	} 
\end{ex}

% nguồn: Phan1_Lop_12_Hinh_hoc_Chuong_5
\begin{ex}%[2H5H1-3]%[ma:0b638cf3]
	Trong không gian $Oxyz$,  cho ba điểm $A(2;0;0), B(0;-3;0), C(0;0;4)$. Phương trình mặt phẳng $(ABC)$~là
	 \choice
		{\True  $\dfrac{x}{2}-\dfrac{y}{3}+\dfrac{z}{4}=1$}
		{$\dfrac{x}{2}+\dfrac{y}{3}+\dfrac{z}{4}=1$} 
		 {$2x-3y+4z-1=0$}
		 {$2x+3y+4z-1=0$}
	 \loigiai{
		Ta có phương trình mặt phẳng qua đoạn  chắn $(ABC)$
		$$\dfrac{x}{2}+\dfrac{y}{-3}+\dfrac{z}{4}=1.$$ 
	 }
\end{ex}

% nguồn: Phan1_Lop_12_Hinh_hoc_Chuong_5
\begin{ex}%[2H5H1-3]%[ma:ea528e6b]
	Trong không gian $Oxyz$, cho điểm $A(1;-1;2)$ và mặt phẳng $(P)\colon 2x-y+z+1=0$. Viết phương trình mặt phẳng $(Q)$  đi qua điểm $A$ và song song với $(P)$.
	\choice 
		{$(Q)\colon 2x+y-z+1=0$} 
		{\True $(Q)\colon  2x-y+z-5=0$}
		 {$(Q)\colon 2x-y+z=0$}
		{$(Q)\colon  x+y+z-2=0$}
	 \loigiai{
		Mặt phẳng $(Q)$ đi qua điểm $A$ và song song với  $(P)$ là
		 $$2(x-1)-1(y+1)+1(z-2)=0\Leftrightarrow 2x-y+z-5=0.$$
	 }
\end{ex}

% nguồn: Phan1_Lop_12_Hinh_hoc_Chuong_5
\begin{ex}%[2H5H1-3]%[ma:68729c10]
	Trong không gian với hệ tọa độ $Oxyz$, phương trình mặt phẳng $(P)$ đi qua điểm $B(2;1;-3)$, đồng thời  vuông góc với hai mặt phẳng $(Q)\colon x + y + 3z - 5 = 0$, $(R)\colon 2x - y + z - 1 = 0$ là
	\choice 
		{$4x + 5y  - 3z + 22 = 0$}
		{$4x - 5y - 3z - 12 = 0$} 
		{$2x + y - 3z - 14  = 0$}
		{\True $4x + 5y -  3z - 22 = 0$}
	 \loigiai{
		Vectơ pháp tuyến  của mặt phẳng $(Q)$ và $(R)$ lần lượt là $\overrightarrow{n_{(Q)}} = (1;1;3)$ và $\overrightarrow{n_{(R)}} = (2;-1;1)$.\\
		$\heva{&(P) \perp (Q) \Rightarrow \overrightarrow{n_{(P)}} \perp \overrightarrow{n_{(Q)}} \\& (P) \perp (R) \Rightarrow \overrightarrow{n_{(P)}} \perp \overrightarrow{n_{(R)}} } \Rightarrow$ vectơ pháp tuyến $\overrightarrow{n_{(P)}} = [\overrightarrow{n_{(Q)}}, \overrightarrow{n_{(R)}}] =  (4;5;-3)$.\\
		Mặt phẳng $(P)$ đi qua điểm $B(2;1;-3)$ và vectơ pháp tuyến $\overrightarrow{n_{(P)}} = (4;5;-3)$  có phương trình
		$$4(x - 2) + 5(y - 1) - 3(z + 3) = 0 \Leftrightarrow 4x  + 5y - 3z - 22 = 0.$$
	 }
\end{ex}

% nguồn: Phan1_Lop_12_Hinh_hoc_Chuong_5
\begin{ex}%[2H5H1-3]%[ma:89c65a70]
	Trong không gian $Oxyz$, phương trình mặt phẳng qua $A(-1; 1;-2)$  và có vectơ pháp tuyến $\overrightarrow{n}=(1;-2;-2)$ là
	 \choice
		 {$-x+y-2z+1=0$}
		 {\True $x-2y-2z-1=0$}
		{$x-2y-2z+7=0$} 
		{$-x+y-2z-1=0$} 
	\loigiai{Phương trình mặt  phẳng qua $A(-1; 1;-2)$ và có vectơ pháp tuyến $\overrightarrow{n}=(1;-2;-2)$ là $1(x+1)-2(y-1)-2(z+2)=0 \Leftrightarrow x-2y-2z-1=0$.}
\end{ex}

% nguồn: Phan1_Lop_12_Hinh_hoc_Chuong_5
\begin{ex}%[2H5H1-3]%[ma:eae28296]
	Trong không gian $Oxyz$,  cho $2$ điểm $M(3;-1;2)$, $N(0;1;-1)$. Mặt phẳng $(P)$ qua $2$ điểm $M$, $N$ và song song với đường thẳng $d:\dfrac{x-1}{2}=\dfrac{y-5}{1}=\dfrac{z+2}{-5}$ có phương trình là $ax+by+cz-2=0$. Tính giá trị biểu thức $T=a+b+c$.
	 \choice
		 {$T=6$}
		{$T=-6$} 
		 {\True $T=5$}
		 {$T=7$}
	 \loigiai{
		Có $\vec{MN}=(-3;2;-3)$ và $\vv{u}=(2;1;-5)$ là một vectơ chỉ phương của $d$  \\
		$\left[\vec{MN},\vec{u}\right]=(-7;-21;-7)$. \\ 
		Mặt phẳng $(P)$ đi qua $M$, $N$ và song song với  $d$ nên $\vec{MN}$ và $\vec{u}$ là cặp vectơ chỉ phương của $(P)$. \\
		Chọn $\vec{n}=-\dfrac{1}{7} \cdot \left[\vec{MN},\vec{u}\right]=(1;3;1)$ là một vectơ pháp tuyến của $(P)$. \\ 
		Phương tình tổng quát của  $(P)$ là \\
		 $1 \cdot (x-3)+3 \cdot (y+1)+1 \cdot (z-2)=0$ hay $x+3y+z-2=0$. \\
		Vậy  $T=1+3+1=5$.
	 }
\end{ex}

% nguồn: Phan1_Lop_12_Hinh_hoc_Chuong_5
\begin{ex}%[2H5H1-3]%[ma:758168ae]
	Trong không gian $Oxyz$, cho hai điểm $A(2; 1;-3)$ và $B(4; 3; 1)$. Mặt phẳng đi qua trung điểm của đoạn thẳng $AB$ và vuông góc  với $AB$ có phương trình là
	 \choice
		{\True $x+y+2z-3=0$} 
		{$x+2y-z+3=0$} 
		 {$x+y+2z+3=0$}
		{$x+2y-z-3=0$} 
	\loigiai{ 
		Gọi $I$ là trung  điểm $AB$ ta có $I(3;2;-1)$.\\
		Mặt phẳng đi qua $I$ và vuông góc với $AB$ nên nhận $\overrightarrow{AB}=(2;2;4)$ là vectơ pháp tuyến có phương trình là 
		$$2(x-3)+2(y-2)+4(z+1)=0\Leftrightarrow x+y+2z-3=0.$$ 
	 }
\end{ex}

% nguồn: Phan1_Lop_12_Hinh_hoc_Chuong_5
\begin{ex}%[2H5H1-3]%[ma:074a7bf9]
	Trong  không gian với hệ tọa độ $Oxyz$, mặt phẳng đi qua điểm $M(1;-1;1)$, $N(2;1;2)$ và song song với trục $Oz$ có phương trình là
	 \choice
		{$x + 2y + z =  0$}
		{\True $2x - y -  3 = 0$}
		 {$x + 2y + z - 6 = 0$}
		 {$2x - y + 5 = 0$}
	 \loigiai{
		Ta có $\overrightarrow{MN} = (1;2;1)$.\\ 
		Trục $Oz$ có vectơ chỉ phương $\overrightarrow{k} =  (0;0;1)$.\\
		 Suy ra
		$[\overrightarrow{MN}, \overrightarrow{k}] = (2;-1;0)$ là vectơ pháp tuyến  của mặt phẳng.\\
		Phương trình mặt phẳng đi qua điểm $M(1;-1;1)$ có vectơ pháp tuyến $(2;-1;0)$ là 
		$$2(x - 1) - 1(y + 1) + 0(z - 1) = 0 \Leftrightarrow 2x - y - 3  = 0.$$
	 }
\end{ex}

% nguồn: Phan1_Lop_12_Hinh_hoc_Chuong_5
\begin{ex}%[2H5H1-3]%[ma:a053e15a]
	Phương trình nào dưới đây là phương trình mặt phẳng đi qua điểm $M(1; 2;-3)$ và có vectơ pháp  tuyến $\overrightarrow{n}=(1;-2; 3)$?
	 \choice
		 {$x-2 y+3 z-12=0$}
		{$x-2  y-3 z+6=0$}
		{\True $x-2  y+3 z+12=0$}
		{$x-2 y-3 z-6=0$} 
	 \loigiai{
		Mặt phẳng đi qua điểm $M(1; 2;-3)$  và có vectơ pháp tuyến $\overrightarrow{n}=(1;-2; 3)$ nên phương trình mặt phẳng là $1(x-1)-2(y-2)+3(z+3)=0\Leftrightarrow x-2y+3z+12=0$.
	} 
\end{ex}

% nguồn: Phan1_Lop_12_Hinh_hoc_Chuong_5
\begin{ex}%[2H5H1-3]%[ma:634ef8a3]
	Trong không  gian với hệ tọa độ $Oxyz$, phương trình mặt phẳng $(P)$ đi qua điểm $N(2;1;-2)$, song song với trục $Oy$ và vuông góc với mặt phẳng $(Q)\colon x - 2y + z - 8 = 0$ là
	 \choice
		{\True $x -  z - 4 = 0$}
		{$x - z + 4 =  0$}
		{$2x + y - 2z -  4 = 0$}
		{$2x + y - 2z +  4 = 0$}
	\loigiai{ 
		Trục $Oy$ có một vectơ  chỉ phương là $\overrightarrow{j} = (0;1;0)$.\\
		Mặt phẳng $(Q)$ một  vectơ pháp tuyến là $\overrightarrow{n_{(Q)}} = (1;-2;1)$.\\
		Mặt phẳng $(P)$ đi qua điểm $N(2;1;-2)$, song song với trục $Oy$  và vuông góc với mặt phẳng $(Q)$ nên có cặp vectơ chỉ phương là $\overrightarrow{j} = (0;1;0)$ và $\overrightarrow{n_{(Q)}} = (1;-2;1)$.\\
		Do  đó $(P)$ có vectơ pháp tuyến là
		$\overrightarrow{n_{(P)}} = \left[\overrightarrow{j}; \overrightarrow{n_{(Q)}}\right] = (1;0;-1)$.\\ 
		Mặt phẳng $(P)$ đi qua điểm $N(2;1;-2)$ và nhận $\overrightarrow{n_{(P)}} = (1;0;-1)$ làm một vectơ pháp tuyến nên có phương trình là $1(x - 2) + 0(y - 1) - 1(z + 2) = 0  \Leftrightarrow x - z - 4 = 0$.
	 }
\end{ex}

% nguồn: Phan1_Lop_12_Hinh_hoc_Chuong_5
\begin{ex}%[2H5H1-3]%[ma:c0b62034]
	Trong không gian với hệ tọa độ $Oxyz$, phương trình mặt phẳng $(P)$ đi qua $3$  điểm $M(-2;0;0)$, $N(0;-1;0)$ và $P(0;0;3)$ là
	 \choice
		{$3x + 6y - 2z  - 6 = 0$}
		{$2x +  y - 3z - 1 = 0$}
		 {\True $3x + 6y - 2z + 6 = 0$}
		{$3x -  6y - 2z = 0$}
	 \loigiai{
		Áp dụng phương trình mặt phẳng theo đoạn chắn, ta có  phương trình mặt $(MNP)$ là
		$$ \dfrac{x}{-2} + \dfrac{y}{-1} + \dfrac{z}{3}  = 1\Leftrightarrow 3x + 6y - 2z + 6 = 0.$$
	 }
\end{ex}

% nguồn: Phan1_Lop_12_Hinh_hoc_Chuong_5
\begin{ex}%[2H5H1-3]%[ma:fb6581a3]
	Trong không gian $Oxyz$, cho điểm $A(1;2;3)$. Mặt phẳng chứa điểm $A$ và trục $Oz$ có  phương trình là
	 \choice
		{\True  $2x-y=0$}
		{$x+y-z=0$} 
		{$3y-2z=0$} 
		 {$3x-z=0$}
	 \loigiai{
		Mặt phẳng  chứa trục $Oz$ nên nhận một véc-tơ chỉ phương là $\overrightarrow{k}=(0;0;1)$.\\
		Mặt phẳng lại đi qua điểm $A(1;2;3)$ nên véc-tơ  $\overrightarrow{OA}=(1;2;3)$ cũng là véc-tơ chỉ phương của mặt phẳng.\\
		Vậy cặp  véc-tơ chỉ phương là $\overrightarrow{k}$ và $\overrightarrow{OA}$.\\
		Véc-tơ  pháp tuyến của mặt phẳng là
		$\overrightarrow{n}=\left[\overrightarrow{k}, \overrightarrow{OA}\right]  =(-2;1;0)=-(2;-1;0)$.\\
		Phương  trình mặt phẳng có dạng $(P)\colon 2x-y+d=0$.\\
		Thay  điểm $O(0;0;0)$ vào phương trình $(P)$ ta được $2\cdot0-0+d=0 \Rightarrow d=0$.\\
		Vậy phương trình mặt phẳng là  $2x-y=0$.
	} 
\end{ex}

% nguồn: Phan1_Lop_12_Hinh_hoc_Chuong_5
\begin{ex}%[2H5H1-3]%[ma:1888f3b5]
	Trong không gian $(Oxyz)$, mặt phẳng  $(P)$ đi qua điểm $A(1;2;3)$ và song song với mặt phẳng $(Q)\colon 2x+3y-4z-5=0$ có phương trình là
	\choice 
		{$2x+3y+4z-14=0$} 
		{$2x-3y-4z+6=0$} 
		{$2x+3y-4z-4=0$} 
		 {\True $2x+3y-4z+4=0$}
	\loigiai{ 
		Mặt phẳng $(Q)\colon 2 x+3 y-4  z-5=0$ có véc-tơ pháp tuyến là $\overrightarrow{n}_{Q}=(2;3;-4)$.\\
		Mặt phẳng $(P)$  song song với mặt phẳng $(Q)$ nên $ \overrightarrow{n}_{Q}=\overrightarrow{n}_{P}=(2;3;-4) $.\\
		Mặt phẳng $ (P) $ đi qua $A(1;2;3)$.\\ 
		Khi đó phương trình mặt phẳng $(P)$  là
		 $$2(x-1)+3(y-2)-4(z-3)=0 \Leftrightarrow 2x+3y-4z+4=0.$$
	} 
\end{ex}

% nguồn: Phan1_Lop_12_Hinh_hoc_Chuong_5
\begin{ex}%[2H5H1-3]%[ma:e1878114]
	Trong không gian $Oxyz$, mặt phẳng  $(\alpha)$ đi qua ba điểm $A(2;3;-1)$, $B(1;0;1)$, $C(0;2;2)$ có phương trình là
	\choice 
		{$7x - y  + 5z - 6 = 0$}
		{\True $7x + y  + 5z - 12 = 0$}
		{$7x - y - 5z - 16 =  0$}
		{$7x +  y - 5z - 22 = 0$}
	\loigiai{ 
		Ta có $\overrightarrow{AB} = (-1;-3;2)$, $\overrightarrow{AC} = (-2;-1;3)$.\\ 
		Suy ra $\overrightarrow{n} = \left[\overrightarrow{AB}; \overrightarrow{AC}\right] = (-7;-1;-5)$ là vectơ  pháp tuyến của mặt phẳng $(\alpha)$.\\
		Vậy phương trình mặt phẳng $(\alpha)$ là $$-7(x - 2) - (y - 3) -5(z + 1) = 0 \Leftrightarrow 7x + y +  5z - 12 = 0.$$
	} 
\end{ex}

% nguồn: Phan1_Lop_12_Hinh_hoc_Chuong_5
\begin{ex}%[2H5H1-3]%[ma:2dc5d4a9]
	Trong hệ tọa độ $Oxyz$, viết phương trình mặt phẳng trung trực của đoạn $A(1;1;-1)$,  $B(5;2;1)$.
	 \choice
		 {$6x+3y-27=0$}
		{\True $8x+2y+4z-27=0$} 
		 {$8x+2y+4z+27=0$}
		 {$4x+y+2z-3=0$}
	\loigiai{ 
		Ta có $\overrightarrow{AB} = (4;1;2)$.  \\
		Gọi $I$ là trung điểm của $AB \Rightarrow I \left( 3; \dfrac{3}{2}; 0 \right)$.\\ 
		Phương trình mặt phẳng trung trực của đoạn thẳng $AB$ đi qua trung điểm $I\left(3; \dfrac{3}{2};0\right)$ của  $AB$ và nhận $\overrightarrow{AB} = (4;1;2)$ làm vectơ pháp tuyến có dạng
		 \[
		4(x-3) + \left(y - \dfrac{3}{2}  \right) + 2z = 0 \Leftrightarrow 8x+2y+4z-27=0.
	\]} 
\end{ex}

% nguồn: Phan1_Lop_12_Hinh_hoc_Chuong_5
\begin{ex}%[2H5H1-3]%[ma:26a139a4]
	Trong không gian $Oxyz$, cho điểm $A(1;2;-1)$ và mặt phẳng $(P)\colon x+2y+z=0$. Mặt  phẳng $(Q)$ qua $A$ và song song với mặt phẳng $(P)$ có phương trình là
	\choice 
		 {$x+2y+z-1=0$}
		 {$x+2y+z+4=0$}
		 {$x+2y+z-6=0$}
		 {\True $x+2y+z-4=0$}
	 \loigiai{
		Mặt phẳng $(Q)$ song song với mặt  phẳng $(P)$ nên có phương trình $$(Q)\colon x+2y+z+c=0\, (c\ne 0).$$
		Ta có $A(1;2;-1)\in (Q)\Rightarrow 1+2\cdot 2-1+c=0\Rightarrow c=-4$.\\ 
		Vậy phương trình mặt phẳng  $(Q)\colon x+2y+z-4=0$.
	 }
\end{ex}

% nguồn: Phan1_Lop_12_Hinh_hoc_Chuong_5
\begin{ex}%[2H5H1-3]%[ma:0a19e442]
	Trong không gian với hệ tọa độ $Oxyz$, cho hai điểm $A(-1;0;1)$ và $B(2;1;0)$. Phương trình mặt phẳng  $(P)$ đi qua $A$ và vuông góc với $AB$ là
	 \choice
		{\True $3x + y  - z + 4 = 0$}
		{$3x + y - z = 0$} 
		{$3x + y  - z = 0$}
		{$2x  + y - z + 1 = 0$}
	\loigiai{ 
		Do mặt phẳng $(P)$ vuông góc $AB$ nên chọn vectơ pháp  tuyến $\overrightarrow{n_{(P)}} = \overrightarrow{AB} = (3;1;-1)$.\\
		Suy ra $(P)\colon 3(x + 1) + y - (z - 1) = 0 \Leftrightarrow (P)\colon 3x + y -  z + 4 = 0$.
	 }
\end{ex}

% nguồn: Phan1_Lop_12_Hinh_hoc_Chuong_5
\begin{ex}%[2H5H1-3]%[ma:7e918581]
	 Trong không gian với hệ tọa độ $Oxyz$, phương trình mặt phẳng $(P)$ đi qua $3$ điểm $A(-1;1;3)$, $B(2;-1;3)$, $C(2;2;-1)$ là
	 \choice
		{\True $8x + 12y + 9z  - 31 = 0$}
		{$8x + 12y + 9z  + 31 = 0$}
		{$8x - 12y + 9z -  31 = 0$}
		{$8x - 12y + 9z + 31 = 0$} 
	\loigiai{ 
		Mặt phẳng $(ABC)$ nhận $\overrightarrow{AB} = (3;-2;0)$, $\overrightarrow{AC} = (3;1;-4)$ làm cặp vectơ chỉ phương nên có  một vectơ pháp tuyến là $\overrightarrow{n} = [\overrightarrow{AB}, \overrightarrow{AC}] = (8;12;9)$.\\
		Mặt  phẳng $(ABC)$ đi qua điểm $A(-1;1;3)$ và nhận $\overrightarrow{n} = (8;12;9)$ làm một vectơ pháp tuyến nên có phương trình là
		$8(x + 1) + 12(y - 1)  + 9(z - 3) = 0 \Leftrightarrow 8x + 12y + 9z - 31 = 0.$
	} 
\end{ex}

% nguồn: Phan1_Lop_12_Hinh_hoc_Chuong_5
\begin{ex}%[2H5H1-3]%[ma:0a146215]
	[GV:Xuan  Vy Pham]
	Trong không gian $Oxyz$, cho điểm $A(0;-3;2)$ và mặt phẳng $(P) \colon 2x-y+3z+5=0$. Mặt phẳng đi qua $A$ và song song với $(P)$ có phương trình  là
	 \choice
		{$2x-y+3z+9=0$} 
		{$2x+y+3z-3=0$} 
		 {$2x+y+3z+3=0$}
		{\True $2x-y+3z-9=0$} 
	\loigiai{Mặt phẳng $(Q)$ song song với mặt phẳng $(P)\colon 2x-y+3z+5=0$ có  phương trình là $(Q) \colon 2x-y+3z+D=0$ ($D \neq 5$).\\
		Vì $A(0;-3;2) \in (Q) \Rightarrow 2 \cdot 0 -(-3)+3 \cdot 2 +D=0  \Rightarrow D=-9$.\\
	 Vậy $(Q) \colon 2x-y+3z-9=0$.}
\end{ex}

% nguồn: Phan1_Lop_12_Hinh_hoc_Chuong_5
\begin{ex}%[2H5H1-3]%[ma:3760ceb5]
	Trong không gian $Oxyz$, cho hai  điểm $A(1;-2;3)$ và $B(3;1;1)$. Đường thẳng $AB$ có phương trình là
	\choice 
		{$\dfrac{x-4}{1}=\dfrac{y+1}{-2}=\dfrac{z-4}{3}$} 
		{$\dfrac{x-1}{4}=\dfrac{y+2}{-1}=\dfrac{z-3}{4}$} 
		{\True  $\dfrac{x-1}{2}=\dfrac{y+2}{3}=\dfrac{z-3}{-2}$}
		 {$\dfrac{x-2}{1}=\dfrac{y-3}{-2}=\dfrac{z+2}{3}$}
	\loigiai{Ta có $\overrightarrow{AB}=(3-1; 1-(-2);  1-3)=(2; 3; -2)$.
		Phương trình chính tắc của đường thẳng $AB$ đi qua $A(1;-2;3)$ và  nhận $\overrightarrow{AB}=(2;3;-2)$ làm VTCP là
		 \[\dfrac{x-1}{2}=\dfrac{y+2}{3}=\dfrac{z-3}{-2}.\]
	} 
\end{ex}

% nguồn: Phan1_Lop_12_Hinh_hoc_Chuong_5
\begin{ex}%[2H5H1-3]%[ma:3c68628d]
	Trong không gian $Oxyz$, cho  ba điểm $A(0;2;1)$, $B(3;0;1)$, $C(1;0;0)$. Phương trình mặt phẳng $(ABC)$ là
	 \choice
		 {$2x-3y-4z+2=0$}
		 {$4x+6y-8z+2=0$}
		{\True  $2x+3y-4z-2=0$}
		{$2x-3y-4z-2=0$} 
	\loigiai{ 
		Ta  có $\overrightarrow{AB}=(3;-2;0)$, $\overrightarrow{AC}=(1;-2;-1)$.\\
		Véc-tơ pháp tuyến của mặt phẳng $(ABC)$  là $\overrightarrow{n}=\left[\overrightarrow{AB}, \overrightarrow{AC}\right]=(2;3;-4)$.\\
		Suy  ra phương trình mặt phẳng $(ABC)$ có dạng $2x+3y-4z+d=0$.\\
		Thay tọa độ điểm $A(0;2;1)$  vào $(ABC)$ ta được $2\cdot0+3\cdot2-4\cdot1+d=0 \Rightarrow d=-2$.\\
		Vậy phương trình mặt phẳng  là $2x+3y-4z-2=0$.\\
	} 
\end{ex}

% nguồn: Phan1_Lop_12_Hinh_hoc_Chuong_5
\begin{ex}%[2H5H1-3]%[ma:57383b70]
	Trong không gian $Oxyz$, cho ba điểm  ${A}\left(1;2;3\right)$, ${B}\left(1;1;1\right)$ và ${C}\left(3;0;2\right)$. Phương trình mặt phẳng $\left({P}\right)$ đi qua $A$ và vuông góc với đường thẳng $BC$~là
	\choice 
		{\True  $ 2x-y+z-3=0$}
		{$  2x-y+z=0$}
		{$ x+2y+3z-1=0$} 
		{$ x+2y+3z-7=0$} 
	\loigiai{Phương trình  mặt phẳng $\left({P}\right)$ đi qua $A$ và vuông góc với đường thẳng $BC$ nên
		$$(P)\colon \heva{&\text{ Qua } A(1;2;3)\\ & \vec{n}=\overrightarrow{BC} =(2;-1;1) \text{ ~là một VTPT}}\Rightarrow (P)\colon  2x-y+z-3=0. $$
	 }
\end{ex}

% nguồn: Phan1_Lop_12_Hinh_hoc_Chuong_5
\begin{ex}%[2H5H1-3]%[ma:bac0d39a]
	Trong không gian với hệ  tọa độ $Oxyz$, phương trình mặt phẳng $(P)$ đi qua điểm $A(0;1;3)$ và song song với mặt phẳng $(Q)\colon 2x - 3z + 2025 = 0$ là
	\choice 
		{\True $(P)\colon 2x - 3z + 9  = 0$}
		 {$(P)\colon 2x - 3z - 9 = 0$}
		{$(P)\colon 2x - 3z  + 3 = 0$}
		{$(P)\colon 2x  - 3z - 3 = 0$}
	\loigiai{ 
		Ta có $(P)$ qua $A(0;1;3)$  và vectơ pháp tuyến $\overrightarrow{n_{(P)}} = \overrightarrow{n_{(Q)}} = (2;0;-3)$.\\
		Phương trình mặt phẳng $(P)$ là  $2(x - 0) + 0(y - 1) - 3(z - 3) = 0$ $\Leftrightarrow$ $2x - 3z + 9 = 0$.
	 }
\end{ex}

% nguồn: Phan1_Lop_12_Hinh_hoc_Chuong_5
\begin{ex}%[2H5H1-3]%[ma:bf817618]
	Trong không gian với hệ trục $Oxyz$, mặt phẳng $(P)$ đi qua điểm $A(-1; -2; 2)$ và chứa trục $Oz$, có một vectơ pháp  tuyến $\vec{n} = (1; a; b)$. Giá trị của $a + b$ là
	\choice 
		{$-\dfrac{3}{2}$} 
		{$-1$} 
		 {\True $-\dfrac{1}{2}$}
		 {$-2$}
	 \loigiai{
		Vì mặt phẳng $(P)$ chứa trục $Oz$ nên nhận vectơ chỉ  phương $\vec{u} = (0; 0; 1)$.\\
		 Mặt khác, $\vec{n} = (1; a; b)$ là pháp tuyến của $(P)$ nên ta có
		$$\vec{n} \cdot \vec{u} = 0 \Rightarrow 1 \cdot 0 + a \cdot 0 +  b \cdot 1 = 0 \Rightarrow b = 0.$$
		Vậy $\vec{n} =  (1; a; 0)$.\\
		Ta có  mặt phẳng $(P)$ đi qua điểm $A(-1; -2; 2)$ và có pháp tuyến $\vec{n} = (1; a; 0)$ nên phương trình mặt phẳng là
		$$1(x +  1) + a(y + 2) + 0(z - 2) = 0 \Rightarrow x + 1 + a(y + 2) = 0.$$
		Để mặt phẳng chứa trục $Oz$  thì điểm $O(0; 0; 0)$ thuộc mặt phẳng. Thay tọa độ $O$ vào
		$$0 + 1 + a(0 + 2) = 0 \Rightarrow 1 + 2a =  0 \Rightarrow a = -\dfrac{1}{2}.$$
		Vì $b = 0$ nên $a + b = -\dfrac{1}{2} + 0  = -\dfrac{1}{2}$.
	 }
\end{ex}

% nguồn: Phan1_Lop_12_Hinh_hoc_Chuong_5
\begin{ex}%[2H5H1-3]%[ma:9774bc34]
	Trong không gian $ Oxyz$, mặt phẳng đi qua ba  điểm $ A\left(1;1;2\right)$, $ B\left(6;0;0\right)$, $ C\left(2;4;0\right)$ có phương trình~là
	 \choice
		{\True $ x+y+2z-6=0$} 
		{$  x+y+z-4=0$}
		{$  x+3y+z-6=0$}
		{$  x+y+z+4=0$}
	\loigiai{Mặt phẳng $(ABC)$ có cặp véc-tơ chỉ phương~là $\overrightarrow{AB}=(5;-1;-2)$ và $\overrightarrow{AC}=(1;3;-2)$ nên có $[\overrightarrow{AB}, \overrightarrow{AC}]= (8;8;16)$ chọn một véc-tơ pháp tuyến~là $\vec{n}=(1;1;2)$.  Vậy $(ABC)\colon x+y+2z-6=0.$
	} 
\end{ex}

% nguồn: Phan1_Lop_12_Hinh_hoc_Chuong_5
\begin{ex}%[2H5V1-3]%[ma:39379a6c]
	Trong không gian $Oxyz$, cho các mặt phẳng $(Q) \colon x+y+3z=0, (R) \colon 2x-y+z=0$ và điểm $B(2;1;-3)$. Mặt phẳng $(P)$ đi qua điểm $B$ đồng thời vuông góc  với hai mặt phẳng $(Q)$ và $(R)$ có phương trình~là
	 \choice
		{\True $4x+5y-3z+22=0$} 
		 {$4x-5y-3z-12=0$}
		 {$2x+y-3z-14=0$}
		 {$4x+5y-3z-22=0$}
	 \loigiai{
		Gọi $(\alpha)$~là mặt phẳng cần tìm.\\ 
		Ta có $(\alpha) \colon \heva{&\text{ Qua  }B(2;1;-3)\\&\vec{n}_{P}=\left[\vec{n}_{Q},\vec{n}_{(R)}\right]=(4;5;-3).}$\\
		 Suy ra $(P) \colon 4(x-2)+5(y-1)-3(z+3)=0 \Leftrightarrow 4x+5y-3z-22=0$
	 .}
\end{ex}

% nguồn: Phan1_Lop_12_Hinh_hoc_Chuong_5
\begin{ex}%[2H5B1-3]%[ma:ece5c933]
	Trong  không gian $Oxyz$, mặt phẳng đi qua điểm $A\left(1;2;3 \right)$ và có vectơ pháp tuyến $\vec{n}=\left(3;-2;-1 \right)$ có phương trình~là
	\choice 
		{$3x-2y-z-4=0$} 
		 {$3x-2y-z=0$}
		{\True  $3x-2y-z+4=0$}
		 {$x+2y+3z+4=0$}
	\loigiai{ 
		Mặt phẳng đi qua  điểm $A\left(1;2;3 \right)$ và có vectơ pháp tuyến $\vec{n}=\left(3;-2;-1 \right)$ có phương trình~là
		$$3(x-1)-2(y-2)-1(z-3)=0 \Leftrightarrow 3x-2y-z+4=0.$$ 
	 }
\end{ex}

% nguồn: Phan1_Lop_12_Hinh_hoc_Chuong_5
\begin{ex}%[2H5B1-3]%[ma:93eaec23]
	 Trong không gian $Oxyz$, phương trình mặt phẳng $\left(Ozx \right)$~là
	\choice 
		 {$x=0$}
		 {\True $y=0$}
		{$z=0$} 
		{$x+y+z=0$} 
	\loigiai{ 
		Phương trình mặt  phẳng $\left(Ozx \right)$~là $y=0$.
	 }
\end{ex}

% nguồn: Phan1_Lop_12_Hinh_hoc_Chuong_5
\begin{ex}%[2H5B1-3]%[ma:f1cf0b38]
	Trong không  gian $Oxyz$, mặt phẳng đi qua điểm ${A}\left(3;0;0\right)$, ${B}\left(0;5;0\right)$ và ${C}\left(0;0;7\right)$ có phương trình~là
	\choice 
		{$\dfrac{y}{3}+\dfrac{x}{5}+\dfrac{z}{7}=1$} 
		 {$\dfrac{z}{3}+\dfrac{y}{5}+\dfrac{x}{7}=1$}
		{$\dfrac{x}{3}+\dfrac{z}{5}+\dfrac{y}{7}=1$} 
		{\True  $\dfrac{x}{3}+\dfrac{y}{5}+\dfrac{z}{7}=1$}
	\loigiai{ 
		Phương  trình mặt phẳng theo đoạn chắn $\big(ABC\big)\colon \dfrac{x}{3}+\dfrac{y}{5}+\dfrac{z}{7}=1.$
	} 
\end{ex}

% nguồn: Phan1_Lop_12_Hinh_hoc_Chuong_5
\begin{ex}%[2H5B1-3]%[ma:16b1dcc7]
	Trong không gian $Oxyz$, cho mặt phẳng $\left(P \right)\colon x-2y+3z+4=0$. Phát biểu nào sau đây  đúng?
	 \choice
		{$\vec{n}_1=\left(1;2;3 \right)$~là một vectơ pháp tuyến của mặt  phẳng $\left(P \right)$}
		{\True $\vec{n}_2=\left(1;-2;3 \right)$~là  một vectơ pháp tuyến của mặt phẳng $\left(P \right)$}
		{$\vec{n}_3=\left(1;3;4 \right)$~là một vectơ  pháp tuyến của mặt phẳng $\left(P \right)$}
		{$\vec{n}_4=\left(-2;3;4 \right)$~là một vectơ  pháp tuyến của mặt phẳng $\left(P \right)$}
	 \loigiai{
		Mặt  phẳng $\left(P \right)\colon x-2y+3z+4=0$ có một vectơ pháp tuyến~là $\vec{n}=\left(1;-2;3 \right)$.
	 }
\end{ex}

% nguồn: Phan1_Lop_12_Hinh_hoc_Chuong_5
\begin{ex}%[2H5B1-3]%[ma:7ae346c2]
	Trong không gian $Oxyz$, mặt phẳng đi qua điểm $M\left(7;8;9 \right)$  vuông góc với $Ox$ có phương trình~là
	 \choice
		 {\True $x=7$}
		{$y=8$} 
		{$z=9$} 
		{$\left(x-7 \right)+\left(y-8 \right)+\left(z-9 \right)=0$} 
	\loigiai{ 
		Mặt phẳng vuông góc với $Ox$ nên nhận $\vec{i}=(1;0;0)$~làm vectơ pháp tuyến.\\ 
		Khi đó mặt phẳng có phương trình~là  $1(x-7)=0 \Leftrightarrow x=7.$
	} 
\end{ex}

% nguồn: Phan1_Lop_12_Hinh_hoc_Chuong_5
\begin{ex}%[2H5B1-3]%[ma:811b6395]
	Trong không gian $Oxyz$, phương trình mặt phẳng $(Oxy)$~là 
	\choice 
		{$x=0$} 
		 {$y=0$}
		{\True  $z=0$}
		{$x+y+z=0$} 
	 \loigiai{
		Phương trình  mặt phẳng $(Oxy)$~là $z=0$.
	} 
\end{ex}


%%%%%%%%%%%%%%%% ĐÚNG/SAI %%%%%%%%%%%%%%%%

% nguồn: Phan1_Lop_12_Hinh_hoc_Chuong_5
\begin{ex}%[2H5N1-3]%[ma:48d7bccc]
	Trong không gian $Oxyz$, cho ba điểm $A(1;  1; 1)$, $B(1; -2; 3)$, $C(2; -1; 2)$.
	 \choiceTF
		{Ba điểm $A, B, C$ đã cho  thẳng hàng}
		{Có vô số mặt phẳng đi  qua ba điểm $A, B, C$ đã cho}
		 {\True Mặt phẳng $(ABC)$ có một vectơ pháp tuyến $\overrightarrow{n} = (1; 2; 3)$}
		{\True Mặt phẳng $(ABC)$ có phương trình là $x  + 2y + 3z - 6 = 0$}
	 \loigiai{
		\begin{itemchoice}
			\itemch $\overrightarrow{AB} = (0; -3; 2)$, $\overrightarrow{AC} = (1; -2; 1)$, $\dfrac{0}{1}  \ne \dfrac{-3}{-2}$ suy ra không cùng phương nên ba điểm không thẳng hàng.
			\itemch Vì ba điểm không thẳng hàng nên chỉ  có duy nhất một mặt phẳng đi qua chúng.
			\itemch $\overrightarrow{n} = [\overrightarrow{AB}, \overrightarrow{AC}] = (1; 2; 3)$. 
			\itemch Mặt phẳng đi  qua $A(1;1;1)$ và có pháp tuyến $\overrightarrow{n} = (1;2;3)$ là\\
			$$1(x - 1) + 2(y - 1) + 3(z - 1)  = 0 \Rightarrow x + 2y + 3z - 6 = 0.$$
		\end{itemchoice}
	} 
\end{ex}

% nguồn: Phan1_Lop_12_Hinh_hoc_Chuong_5
\begin{ex}%[2H5H1-3]%[ma:e703fd81]
	Trong không gian $Oxyz$, cho các điểm $A(1; 1;  3)$, $B(-1; 3; 2)$, $C(-1; 2; 3)$.
	\choiceTF 
		{Vectơ $\overrightarrow{AB}$ và $\overrightarrow{AC}$  cùng phương}
		{\True Điểm $G\left(\dfrac{-1}{3}; 2; \dfrac{8}{3}\right)$ là trọng tâm tam giác  $ABC$}
		{Mặt phẳng $(ABC)$ có phương trình là $x + ay + bz  + c = 0$. Khi đó $a + b + c = 4$}
		{Điểm $M(1; 2; 3)$  thuộc mặt phẳng $(ABC)$}
	 \loigiai{
		\begin{itemchoice}
			\itemch $\overrightarrow{AB} = (-2; 2; -1)$, $\overrightarrow{AC} = (-2; 1; 0)$, kiểm tra tỷ lệ $\dfrac{-2}{-2} \ne \dfrac{1}{2} \ne \dfrac{0}{-1}$ suy ra không cùng phương. 
			\itemch Đúng. 
			\itemch $\overrightarrow{AB} = (-2; 2; -1)$, $\overrightarrow{AC} = (-2; 1; 0)$ $\Rightarrow [\overrightarrow{AB}, \overrightarrow{AC}] = (1; 2; 2)$.\\ 
			Mặt phẳng $(ABC)$ qua $A$  và có pháp tuyến $(1;2;2)$ nên ta có phương trình
			$$1(x - 1) + 2(y - 1) + 2(z - 3) = 0 \Rightarrow x + 2y  + 2z - 9 = 0.$$
			Khi đó $a = 2$, $b = 2$, $c = -9$ $\Rightarrow a + b +  c = -5$.
			\itemch Thay $M(1;2;3)$ vào $x +  2y + 2z - 9 = 0$ là $1 + 4 + 6 - 9 = 2 \ne 0$.
		\end{itemchoice}
	} 
\end{ex}

% nguồn: Phan1_Lop_12_Hinh_hoc_Chuong_5
\begin{ex}%[2H5H1-3]%[ma:18b0f07a]
	Cho các điểm $A(1; -2; 0)$, $B(2; -1; 1)$, $C(1; 1; 2)$. 
	\choiceTF 
		{Phương trình mặt phẳng  $(ABC)$ là $x+2y-3z-3=0$}
		{\True Phương trình mặt phẳng $(\alpha)$  qua $A$ và vuông góc với $BC$ là $x-2y-z-5=0$}
		{\True Phương trình  mặt phẳng trung trực $(\beta)$ của đoạn $AC$ là $6y+4z-1=0$}
		 {Phương trình mặt phẳng $(\gamma)$ chứa trục $Ox$ và điểm $C$ là $2y+z=0$}
	 \loigiai{
		\begin{itemchoice}
			\itemch Ta có $\overrightarrow{AB} = (1; 1; 1)$,  $\overrightarrow{AC} = (0; 3; 2)$ suy ra $\overrightarrow{n} = [\overrightarrow{AB}, \overrightarrow{AC}] = (-1; -2; 3)$ là một vectơ pháp tuyến của mặt phẳng.\\
			Phương trình mặt phẳng đi qua $A(1; -2; 0)$  là
			$$-1(x - 1) - 2(y + 2) +  3z = 0 \Rightarrow x + 2y - 3z + 3 = 0.$$
			Phương trình mặt  phẳng đã cho là $x + 2y - 3z - 3 = 0$.
			\itemch Ta có $\overrightarrow{BC} = (-1; 2; 1)$ là vectơ pháp tuyến của mặt  phẳng.\\
			Mặt phẳng đi qua $A(1; -2; 0)$ và vuông góc với $BC$ có phương  trình
			$$-1(x -  1) + 2(y + 2) + 1z = 0 \Rightarrow x - 2y - z - 5 = 0.$$
			\itemch Trung điểm $AC$ là  $M\left(1; \dfrac{-1}{2}; 1\right)$, $\overrightarrow{AC} = (0; 3; 2)$ là vectơ pháp tuyến.\\
			Phương trình mặt phẳng  trung trực là\\
			$$0(x - 1)  + 3\left(y + \dfrac{1}{2}\right) + 2(z - 1) = 0 \Rightarrow 6y + 4z - 1 = 0.$$
			\itemch Ta có $\overrightarrow{i} = (1; 0; 0)$  và $\overrightarrow{OC} = (1; 1; 2)$.\\
			Pháp tuyến mặt phẳng chứa $Ox$  và $C$ là $[\overrightarrow{i}, \overrightarrow{OC}] = (0; -2; 1)$.\\
			Phương trình mặt phẳng là $-2y + z  = 0 \Leftrightarrow 2y - z = 0$.\\
			Phương trình mặt phẳng  đã cho là $2y + z = 0$.
		\end{itemchoice}
	 }
\end{ex}

% nguồn: Phan1_Lop_12_Hinh_hoc_Chuong_5
\begin{ex}%[2H5H1-3]%[ma:f30a361e]
	Trong không gian $Oxyz$, cho hai điểm $A(1;2;-2)$, $B(2;-2;1)$ và  mặt phẳng $(Q)\colon x+3y+z+3=0$.
	 \choiceTF
		{\True Điểm $B$ không thuộc mặt phẳng  $(Q)$}
		{\True Mặt phẳng  $Q$ có một vectơ pháp tuyến là $\vec{n}=(1;3;1)$}
		{\True Mặt phẳng $P$ đi qua hai  điểm $A$, $B$ và vuông góc với mặt phẳng $(Q)$ có phương trình là $-13x+2y+7z+23=0$}
		{Góc giữa đường thẳng $AB$ và mặt phẳng $(Q)$ (đơn vị: độ, làm tròn  đến hàng phần chục) là $61{,}8^\circ$}
	 \loigiai{
		\begin{itemchoice}
			\itemch Thay tọa độ điểm $B$ vào $(Q)$  ta được $2+3\cdot(-2)+1+3=0$ nên $B\in(Q)$.
			\itemch Vectơ pháp tuyến của  $(Q)$ là $\vec{n}_{P}=(1;3;1)$.
			\itemch  Ta có $\overrightarrow{AB}=(1;-4;3)$ và $\vec{n}_{Q}=(1;3;1)$.\\
			Ta có  vectơ pháp tuyến của mặt phẳng $(P)$ $\vec{n}_{P}=[\overrightarrow{AB},\vec{n}_{Q}]=(-13;2;7)$.\\
			 Phương trình mặt phẳng $(P)$ đi qua $A(1;2;-2)$ có vectơ pháp tuyến $\vec{n}_{P}=(-13;2;7)$ có dạng $-13(x-1)+2(y-2)+7(z+2)=0\Leftrightarrow -13x+2y+7z+23=0$.
			 \itemch $\sin\left(\overrightarrow{AB}, \vec{n}_{Q}\right)=\dfrac{|1-12+3|}{\sqrt{26}\cdot\sqrt{11}}=\dfrac{4\sqrt{286}}{143}$.\\
			Suy  ra $(AB, (Q))\approx28{,}2^\circ$.
		\end{itemchoice}
	} 
\end{ex}

% nguồn: Phan1_Lop_12_Hinh_hoc_Chuong_5
\begin{ex}%[2H5H1-3]%[ma:bdac7331]
	Trong không gian $Oxyz$, cho ba điểm $A\left( 1;2;0 \right)$, $B\left( 1;0;2 \right)$, $C\left( 2;1;3 \right)$  và mặt phẳng \break $\left( P \right) \colon x-y+2z+7=0$.
	\choiceTF 
		{Mặt phẳng $\left( ABC \right)$ có một vectơ pháp tuyến là $\left( 2;1;1 \right)$} 
		{\True Mặt phẳng $\left( ABC \right)$ đi qua điểm  $M\left( 3;1;5 \right)$}
		{Mặt phẳng $\left( ABC \right)$ vuông góc với mặt phẳng $\left( P  \right)$}
		{Khoảng cách từ điểm $A$ đến mặt phẳng $\left(  P \right)$ bằng $6$}
	\loigiai{ 
		\begin{itemchoice}
			\itemch \textbf{Sai}. Ta  có $$\overrightarrow{AB}=\left( 0;-2;2 \right),\,\overrightarrow{BC}=\left( 1;1;1 \right)\Rightarrow \left[\overrightarrow{AB},\overrightarrow{BC} \right]=\left( -4;2;2 \right)=2\left( -2;1;1 \right)=2\cdot\overrightarrow{n}.$$
			Mặt phẳng $\left( ABC \right)$ có một vectơ pháp  tuyến là $\overrightarrow{n}=\left( -2;1;1 \right).$
			 \itemch \textbf{Đúng}. Phương trình mặt phẳng $$\left( ABC \right) \colon -2\left( x-1 \right)+1\left( y-2 \right)+1\left( z-0 \right)=0\Leftrightarrow -2x+y+z=0.$$
			Toạ độ điểm $M\left( 3;1;5 \right)$ thõa  mãn phương trình mặt phẳng $\left( ABC \right)$.
			\itemch \textbf{Sai}. Ta có vectơ pháp tuyến của mặt  phẳng $\left( P \right) \colon \overrightarrow{n_(P)}=\left( 1;-1;2 \right)$.\\
			Khi đó $\overrightarrow{n}_{(P)}\cdot \overrightarrow{n}=-2\cdot 1-1  \cdot 1+1 \cdot 2=-1\ne 0$.\\
			Mặt phẳng $\left(  ABC \right)$ không vuông góc với mặt phẳng $\left( P \right)$.
			\itemch \textbf{Sai}. Ta có $\mathrm{d}\left( A;\left(  P \right) \right)=\dfrac{\left| 1-2+2\cdot 0+7 \right|}{\sqrt{1+1+4}}=\sqrt{6}.$
		\end{itemchoice}
	 }
\end{ex}

% nguồn: Phan1_Lop_12_Hinh_hoc_Chuong_5
\begin{ex}%[2H5H1-3]%[ma:4db39787]
	Cho  các điểm $A\left(1;-2;0\right)$; $B\left(2;-1;1\right)$; $C\left(1;1;2\right)$.
	\choiceTF 
		{Phương trình mặt phẳng $\left(ABC\right)$ là  $x+2y-3z-3=0$}
		{\True Phương trình mặt phẳng $\left(\alpha\right)$ qua $A$ và vuông góc  với $BC$ là $x-2y-z-5=0$}
		{\True Phương trình mặt  phẳng trung trực $\left(\beta\right)$ của đoạn $AC$ là $6y+4z-1=0$}
		{Phương trình mặt phẳng  $\left(\gamma\right)$ chứa trục $Ox$ và điểm $C$ là $2y+z=0$}
	\loigiai{ 
		\begin{itemchoice}
			\itemch Ta  có $\overrightarrow{AB}=\left(1;1;1\right)$; $\overrightarrow{AC}=\left(0;3;2\right)$.\\
			Véc-tơ pháp tuyến  của $\left(ABC\right)$ là $\overrightarrow{n}=\left[\overrightarrow{AB},\overrightarrow{AC}\right]=\left(-1;-2;3\right)$.\\
			Phương trình  mặt phẳng $\left(ABC\right)$ là
			\begin{eqnarray*}
				 -1\cdot\left(x-1\right)-2\cdot\left(y+2\right)+3z=0\Leftrightarrow x+2y-3z+3=0.
			\end{eqnarray*}
			\itemch Véc-tơ  pháp tuyến của $\left(\alpha\right)$ là $\overrightarrow{n}=\overrightarrow{BC}=\left(-1;2;1\right)$.\\
			Phương trình mặt phẳng $\left(\alpha\right)$  là $-1\cdot\left(x-1\right)+2\cdot\left(y+2\right)+1z=0$ hay $x-2y-z-5=0$.
			\itemch Ta có trung điểm của đoạn $AC$  là $M\left(1;\dfrac{-1}{2};1\right)$.\\
			Véc-tơ pháp tuyến của $\left(\beta\right)$ là $\overrightarrow{n}=\overrightarrow{AC}=\left(0;3;2\right)$.\\ 
			Phương trình mặt phẳng  $\left(\beta\right)$ là
			\begin{eqnarray*}
				&&0\cdot\left(x-1\right)+3\cdot\left(y+\dfrac{1}{2}\right)+2\cdot\left(z-1\right)=0\Leftrightarrow  6y+4z-1=0.
			\end{eqnarray*}
			\itemch Ta  có $\overrightarrow{i}=\left(1;0;0\right)$; $\overrightarrow{OC}=\left(1;1;2\right)$.\\
			Véc-tơ pháp tuyến của  $\left(\gamma\right)$ là $\overrightarrow{n}=\left[\overrightarrow{i},\overrightarrow{OC}\right]=\left(0;-2;1\right)$.\\
			Phương trình mặt phẳng $\left(ABC\right)$ là  $0\cdot x-2y+1z=0$ hay $2y-z=0$.
		\end{itemchoice}
	 }
\end{ex}

% nguồn: Phan1_Lop_12_Hinh_hoc_Chuong_5
\begin{ex}%[2H5H1-3]%[ma:d6836fb8]
	Trong không gian $(Oxyz)$, gọi $(\alpha)\colon ax+by+cz-60=0$, $(a,b,c \in \mathbb{R})$ là mặt phẳng thay đổi thỏa mãn luôn đi qua điểm $M(2;4;5)$ và cắt ba tia  $Ox$, $Oy$, $Oz$ tương ứng tại ba điểm $A$, $B$, $C$ khác $O$. Khi đó
	 \choiceTF
		 {\True $2a+4b+5c-60=0$}
		{$A\left(\dfrac{60}{a};0;0\right)$ với $a<0$} 
		{$B\left(0;\dfrac{60}{b};  0\right)$ với $b<0$}
		{\True Thể  tích của khối tứ diện $OABC$ là $V=\dfrac{36\,000}{abc}$}
	\loigiai{ 
		\begin{itemchoice}
			 \itemch \textbf{Đúng}.\\
			Do mặt phẳng $(\alpha)$ đi qua $M(2;4;5)$  ta có $2a+4b+5c-60=0$.
			\itemch \textbf{Sai}.\\ 
			Vì $A$, $B$, $C$ khác $O$  lần lượt thuộc tia $Ox$, $Oy$, $Oz$ nên $x_A>0$, $y_B>0$, $z_C>0$.\\
			Ta có $ax+by+cz-60=0 \Leftrightarrow \dfrac{x}{\tfrac{60}{a}}+\dfrac{y}{\tfrac{60}{b}}+\dfrac{z}{\tfrac{60}{c}}=1 \Rightarrow\heva{&(\alpha) \cap  Ox=A\left(\dfrac{60}{a};0;0\right),\\
			&(\alpha)\cap  Oy=B\left(0;\dfrac{60}{b};0\right),\\
			&(\alpha)\cap Oz=C\left(0;0;\dfrac{60}{c}\right)}$ với  $(a>0, b>0, c>0)$.
			 \itemch \textbf{Sai}.\\
			Ta có $B\left(0;\dfrac{60}{b};0\right)$ với  $b>0$.
			 \itemch \textbf{Đúng}.\\
			Thể tích khối tứ  diện $OABC$ là $V=\dfrac{1}{6}\cdot OA\colon OB\cdot OC=\dfrac{1}{6}\cdot\dfrac{60}{a} \cdot \dfrac{60}{b}\cdot\dfrac{60}{c}=\dfrac{36\,000}{abc}$.
		\end{itemchoice}
	} 
\end{ex}

% nguồn: Phan1_Lop_12_Hinh_hoc_Chuong_5
\begin{ex}%[2H5H1-3]%[ma:5bd2421b]
	Trong không gian với hệ tọa độ $Oxyz$, cho ba  điểm $A(a; 0; 0)$, $B(0; b; 0)$, $C(0; 0; c)$ với $a$, $b$, $c$ đều là số thực dương.
	 \choiceTF
		{\True Mặt phẳng $(ABC)$ có phương trình  $\dfrac{x}{a} + \dfrac{y}{b} + \dfrac{z}{c} = 1$}
		{Mặt phẳng $(ABC)$ đi qua điểm $G(1; 2; 3)$ là trọng tâm của $\triangle ABC$ có phương trình $6x + 3y + 2z +  18 = 0$}
		{\True Mặt  phẳng $(ABC)$ đi qua điểm $H(1; 1; 1)$ là trực tâm của $\triangle ABC$ có phương trình $x + y + z - 3 = 0$}
		{\True Với $OA, OB, OC$ tạo cấp số cộng có công sai $2$ và mặt phẳng $(\alpha)$ đi qua $M(2; -2; 3)$ thì khoảng cách từ $O$ đến mặt phẳng  là $\dfrac{12}{7}$ suy ra $T = 19$}
	 \loigiai{
		\begin{itemchoice}
			\itemch Phương trình đoạn chắn của mặt phẳng đi qua $A$, $B$, $C$ là $\dfrac{x}{a}  + \dfrac{y}{b} + \dfrac{z}{c} = 1$.
			\itemch G  là trọng tâm \\
			$\Rightarrow \dfrac{a}{3} = 1 \Rightarrow a = 3$, $\dfrac{b}{3} = 2 \Rightarrow b = 6$, $\dfrac{c}{3} = 3 \Rightarrow  c = 9$.\\
			Phương trình mặt phẳng là 
			$$\dfrac{x}{3} + \dfrac{y}{6}  + \dfrac{z}{9} = 1 \Leftrightarrow 6x + 3y + 2z - 18 = 0.$$
			 \itemch Với $H(1;1;1)$ là trực tâm, suy ra $\dfrac{1}{a} + \dfrac{1}{b} + \dfrac{1}{c} = 1$ và $b=c$, $a=c$ $\Rightarrow a=b=c=3$.\\
			Phương trình mặt phẳng $$\dfrac{x}{3} + \dfrac{y}{3} + \dfrac{z}{3} = 1 \Leftrightarrow x + y + z = 3.$$ 
			\itemch Gọi mặt phẳng  $(\alpha)\colon \dfrac{x}{a} + \dfrac{y}{a+2} + \dfrac{z}{a+4} = 1$.\\
			Thế $M(2; -2; 3)$ vào ta  được $\dfrac{2}{a} - \dfrac{2}{a+2} + \dfrac{3}{a+4} = 1 \Rightarrow a = 2$.\\
			Từ đó ta được $A(2;0;0)$,  $B(0;4;0)$, $C(0;0;6)$.\\
			Ta  có $$\dfrac{1}{d^2} = \dfrac{1}{OA^2} + \dfrac{1}{OB^2} + \dfrac{1}{OC^2} = \dfrac{1}{4} + \dfrac{1}{16} + \dfrac{1}{36} = \dfrac{49}{144} \Rightarrow d = \dfrac{12}{7} \Rightarrow T = 12 + 7 = 19.$$
		\end{itemchoice}
	} 
\end{ex}

% nguồn: Phan1_Lop_12_Hinh_hoc_Chuong_5
\begin{ex}%[2H5H1-3]%[ma:ba77d66d]
	 Trong không gian $Oxyz$, cho đường thẳng $d\colon \dfrac{x+1}{2}=\dfrac{y-1}{1}=\dfrac{z+2}{-1}$ và mặt phẳng \break $(P)\colon 2 x-y+3 z-1=0$.
	\choiceTF 
		{\True Vectơ  $\overrightarrow{u}=(2; 1;-1)$ là vectơ chỉ phương của đường thẳng $d$}
		{Gốc tọa độ cách  mặt phẳng $(P)$ một đoạn bằng $10$}
		{\True Đường thẳng $d$ song song mặt  phẳng $(P)$}
		{\True Mặt phẳng $(Q)$ chứa đường thẳng $d$ và vuông góc mặt phẳng $(P)$ có phương trình là $x-4 y-2  z+1=0$}
	 \loigiai{
		\begin{itemchoice}
			\itemch Đường thẳng $d\colon \dfrac{x+1}{2}=\dfrac{y-1}{1}=\dfrac{z+2}{-1}$ có 1 VTCP  là $\overrightarrow{u}_d=(2;1;-1)$.
			\itemch  $\mathrm{d}\left[O;(P)\right]=\dfrac{\left|-1\right|}{\sqrt{4+1+9}}=\dfrac{1}{\sqrt{14}}$.
			\itemch Mặt  phẳng $(P)$ có VTPT là $\overrightarrow{n}_{(P)}=(2;-1;3)\Rightarrow \overrightarrow{u}_d.\overrightarrow{n}_{(P)}=2\cdot2-1\cdot1-1\cdot3=0 \Rightarrow \overrightarrow{u}_d\perp \overrightarrow{n}_{(P)}$.\\
			Thế tọa độ điểm $A(-1;1;-2)\in d$ vào phương  trình mp $(P)$ ta có
			$$2\cdot(-1)-1+3\cdot(-2)-1=-10\neq 0 \Rightarrow A \notin  (P) \Rightarrow d \parallel (P).$$
			\itemch Mặt phẳng $(Q)$ chứa $d$ và vuông góc $(P)$ nên nhận $\overrightarrow{u}_d=(2;1;-1)$ và $\overrightarrow{n}_{(P)}=(2;-1;3)$ làm cặp VTCP.\\ 
			Ta có $\left[\overrightarrow{u}_d,\overrightarrow{n}_{(P)}\right]=(2;-8;-4)  \Rightarrow \overrightarrow{n}_{(Q)}=(1;-4;-2)$ là VTPT của $(Q)$ và $A(-1;1;-2)\in d\subset(Q)$.\\
			 PTTQ $(Q)\colon (x+1)-4(y-1)-2(z+2)=0\Leftrightarrow x-4y-2z+1=0$.
		\end{itemchoice}
	 }
\end{ex}

% nguồn: Phan1_Lop_12_Hinh_hoc_Chuong_5
\begin{ex}%[2H5H1-3]%[ma:77117afc]
	Trong không gian với hệ trục tọa độ $Oxyz$, cho hai điểm $A(1; -2; 1)$, $B(0; 1; -3)$  và mặt phẳng $(P)\colon x-y-3=0$.
	\choiceTF 
		 {\True Điểm $A$ thuộc mặt phẳng $(P)$}
		{Mặt phẳng $(Q)$ song song với mặt phẳng $(P)$,  cách $(P)$ một khoảng bằng $2\sqrt{2}$ và cắt trục $Ox$ tại điểm có hoành độ dương có phương trình $(Q)\colon x - y - 1 = 0$}
		{\True Mặt  phẳng $(P)$ song song với trục $Oz$}
		{Mặt phẳng $(\alpha)$ đi qua $A$, $B$ vuông góc với mặt  phẳng $(P)$ có phương trình \break $2x + 2y + z - 1 = 0$}
	 \loigiai{
		\begin{itemchoice}
			\itemch Thay tọa độ $A$ vào phương trình mặt phẳng $(P)$ thấy thỏa mãn nên  $A\in (P)$.
			\itemch  Vì $(Q)\parallel (P)$ nên $(Q)\colon x-y+m=0$, $m\neq -3$.\\
			Giao  điểm của $(Q)$ với trục $Ox$ là $M(-m;0;0)$, do $M$ có hoành độ dương nên $m<0$.\\
			Do $A(1;-;1)\in  (P)$ nên $\mathrm{d}((Q),(P))=\mathrm{d}(A,(Q))=2\sqrt{2}$.\\
			Khi đó $\dfrac{|1-(-2)+m|}{\sqrt{2}}=2\sqrt{2}\Leftrightarrow  |m+3|=4\Leftrightarrow \hoac{&m=1\\&m=-7.}$\\
			 Vì $m<0$ và $m\neq -3$ nên $m=-7$ (thỏa mãn).
			\itemch Mặt phẳng $(P)$ có  véctơ pháp tuyến là $\overrightarrow{n}_P=(1;-1;0)$ và trục $Oz$ có véctơ chỉ phương là $\overrightarrow{k}=(0;0;1)$.\\
			Vì $\overrightarrow{n}_P\cdot \overrightarrow{k}=0$ và  $O\not\in (P)$ nên $Oz\parallel (P)$.
			\itemch Ta có $\overrightarrow{AB}=(-1;3;-4)$, $\left[\overrightarrow{AB};\overrightarrow{n}_P\right]=(-4;-4;-2)=-2(2;2;1)$.\\ 
			Suy ra $\overrightarrow{n}_{(\alpha)}=(2;2;1)$ là một véctơ pháp tuyến của mặt  phẳng $(\alpha)$.\\
			Phương trình mặt phẳng $(\alpha)$  là $2x+2y+z+1=0$.
		\end{itemchoice}
	 }
\end{ex}

% nguồn: Phan1_Lop_12_Hinh_hoc_Chuong_5
\begin{ex}%[2H5V1-3]%[ma:1af52739]
	Trong không gian với hệ toạ độ $Oxyz$,  cho điểm $M(1;0;1)$ và hai mặt phẳng $(Q)\colon x + y - 3z - 5 = 0$, $(R)\colon x + 2y - z - 1 = 0$.
	 \choiceTF
		{Mặt phẳng $(R)$  đi qua điểm $M$}
		{Khoảng cách từ điểm $M$ đến mặt phẳng  $(Q)$ bằng $ \dfrac{7}{\sqrt{11}}$}
		{\True Mặt phẳng đi qua $M$ và song song với mặt phẳng $(Q)$ có phương  trình là $x + y - 3z - 2 = 0$}
		{\True Mặt phẳng đi qua điểm $A(1;-2;0)$ và vuông góc với hai mặt phẳng $(Q) \colon x + y - 3z - 5 = 0$, $(R)\colon x + 2y - z -  1 = 0$ có phương trình là $5x - 2y + z - 9 = 0$}
	 \loigiai{
		Mặt phẳng $(Q)$ có vectơ pháp tuyến là  $\overrightarrow{n}_1=(1;1;-3)$.\\
		Mặt phẳng $(R)$ có vectơ pháp tuyến là $\overrightarrow{n}_2=(1;2;-1)$ 
		\begin{itemchoice}
			\itemch \textbf{Đúng}. Thay tọa độ điểm $M$ vào phương trình mặt phẳng $(R)$ ta được $1+2\cdot 0-1-1 \neq 0$ nên mặt phẳng $(R)$ không đi qua điểm $M$. 
			\itemch \textbf{Sai}. Ta có $\mathrm{d}(M;(Q))  = \dfrac{|1+0-3-5|}{\sqrt{1^2+1^2+3^2}} = \dfrac{7}{\sqrt{11}}$.
			\itemch \textbf{Đúng}.  Gọi $(\alpha)$ là mặt phẳng cần tìm. Do $(\alpha)$ song song với $(Q)$ nên có vectơ pháp tuyến $\overrightarrow{n}=\overrightarrow{n}_1=(1;1;-3)$.\\
			Mặt phẳng $(\alpha)$ đi qua điểm $M$ và có vectơ pháp  tuyến $\overrightarrow{n}=(1;1;-3)$ có phương trình $(x-1)+(y-0)-3(z-1)=0 \Leftrightarrow x+y-3z+2=0$.
			\itemch \textbf{Đúng}. Gọi mặt phẳng cần tìm  là $(P)$. \\
			Mặt phẳng $(P)$ vuông góc với hai mặt phẳng $(Q),(R)$ nên $(P)$ có cặp  vectơ chỉ phương là $\overrightarrow{n}_1=(1;1;-3)$, $\overrightarrow{n}_2=(1;2;-1)$. Do đó, mặt phẳng $(P)$ có vectơ pháp tuyến là
			\[\overrightarrow{n} = [\overrightarrow{n}_1; \overrightarrow{n}_2] =  (5;-2;1) \]
			\noindent Mặt phẳng $(P)$ đi qua điểm $A(1;-2;0)$ và có vectơ pháp tuyến  $\overrightarrow{n}=(5;-2;1)$ có phương trình
			\[5(x-1) - 2(y+2) + (z-0) =  0 \Leftrightarrow 5x - 2y + z - 9 = 0. \]
		\end{itemchoice}
	} 
\end{ex}

% nguồn: Phan1_Lop_12_Hinh_hoc_Chuong_5
\begin{ex}%[2H5V1-3]%[ma:015cce3a]
	Trong không gian tọa độ $Oxyz$, cho mặt phẳng $(P)\colon  x+2y-2z+3=0$ và đường thẳng $d\colon\dfrac{x-1}{2}=\dfrac{y+1}{-1}=\dfrac{z}{1}$.
	 \choiceTF
		{\True Một véc-tơ pháp tuyến của  mặt phẳng $(P)$ là $\overrightarrow{n}=(2; 4;-4)$}
		{\True Góc giữa $d$ và $(P)$ gần bằng  $15^{\circ}47'35{,}41"$}
		{\True Đường thẳng $d$ vuông góc với  đường thẳng $d'\colon\dfrac{x}{3}=\dfrac{y+1}{2}=\dfrac{z}{-4}$}
		{Phương trình mặt phẳng $(Q)$ chứa đường thẳng $d$ và vuông góc với $(P)$ là  \\$x+y+z+1=0$}
	 \loigiai{
		\begin{itemchoice}
			\itemch \textbf{Đúng}. Mặt phẳng $(P)\colon x+2y-2z+3=0$ có véc-tơ pháp tuyến là  $\overrightarrow{n}_1=(1;2;-2)$.\\
			Suy ra  $\overrightarrow{n}_1=(2;4;-4)=2\left(1;2;-2\right)
			=2\overrightarrow{n}$ cũng là một véc-tơ pháp tuyến của mặt  phẳng $(P)\colon x+2y-2z+3=0$.
			\itemch \textbf{Đúng}. Gọi $\alpha$ là góc giữa  đường thẳng $d$ và mặt phẳng $(P)$.\\
			Mặt phẳng $(P)\colon x+2y-2z+3=0$ có  véc-tơ pháp tuyến là $\overrightarrow{n}=(1;2;-2)$.\\
			Đường thẳng  $d\colon\dfrac{x-1}{2}=\dfrac{y+1}{-1}=\dfrac{z}{1}$ có véc-tơ chỉ phương là $\overrightarrow{u}=(2;-1;1)$.\\
			Ta  có $\sin\alpha=\left|\cos(\overrightarrow{u},\overrightarrow{n})\right|=\dfrac{\left|\overrightarrow{u}\cdot\overrightarrow{n}\right|}{\left|\overrightarrow{u}\right|\cdot\left|\overrightarrow{n}\right|}=\dfrac{\left|2\cdot1-1\cdot2-1\cdot2\right|}{\sqrt{2^2+(-1)^2+1^2}\cdot\sqrt{1^2+2^2+(-2)^2}}
			=\dfrac{\sqrt{6}}{9}$.\\ 
			 Suy ra $\alpha\approx 15^{\circ}47'35{,}41"$.
			\itemch  \textbf{Đúng}.
			Đường thẳng $d\colon\dfrac{x-1}{2}=\dfrac{y+1}{-1}=\dfrac{z}{1}$ có véc-tơ  chỉ phương là $\overrightarrow{u}_d=(2;-1; 1)$.\\
			Đường  thẳng $d'\colon\dfrac{x}{3}=\dfrac{y+1}{2}=\dfrac{z}{-4}$ có véc-tơ chỉ phương là $\overrightarrow{u}_d'=(3; 2;-4)$.\\
			Vì  $\overrightarrow{u}_d\cdot\overrightarrow{u}_{d'}=2\cdot3-1\cdot2-1\cdot4=0$ nên $d\perp d'$.
			\itemch  \textbf{Sai}.\\
			Mặt phẳng $(Q)$ chứa đường thẳng $d$  nên $\overrightarrow{u}_d=(2;-1; 1)$ là véc-tơ chỉ phương của $(Q)$. \\
			Mặt phẳng $(Q)$ vuông góc với mặt  phẳng $(P)$ nên véc-tơ pháp tuyến $\overrightarrow{n}_P=(1; 2;-2)$ của $(P)$ là véc-tơ chỉ phương của $(Q)$.\\
			Vậy véc-tơ pháp tuyến của $(Q)$ là $\overrightarrow{n}_Q=[\overrightarrow{u}_d,\overrightarrow{n}_P]=(0; 5; 5)$.\\ 
			Ta có thể  chọn $\overrightarrow{n}_Q=(0; 1;1)$.\\
			Phương trình mặt phẳng $(Q)$  có dạng $0x+1y+1z+D=0\Rightarrow y+z+D=0$.\\
			Vì $(Q)$ chứa $d$ nên  điểm $M_0(1;-1; 0)$ thuộc $d$ cũng phải thuộc $(Q)$.\\
			Thay tọa  độ $M_0$ vào phương trình $(Q)\colon-1+D=0\Rightarrow D=1$.\\
			 Vậy phương trình mặt phẳng $(Q)$ là $y+z+1=0$.
		\end{itemchoice}
	 }
\end{ex}

% nguồn: Phan1_Lop_12_Hinh_hoc_Chuong_5
\begin{ex}%[2H5V1-3]%[ma:c85c370b]
	Trong không  gian $Oxyz$, cho mặt phẳng $(P): 2 x-3 y+z-1=0$ và điểm $A(1;-2; 0)$.
	\choiceTF 
		{\True Một vectơ pháp tuyến của mặt  phẳng $(P)$ là $\overrightarrow{n}=(2;-3; 1)$}
		{Khoảng cách từ điểm $A$ đến mặt phẳng $(P)$ bằng  $\dfrac{4 \sqrt{14}}{7}$}
		{Mặt phẳng $(P)$ đi  qua điểm $B(1; 1; 1)$}
		{\True Phương trình mặt phẳng chứa  điểm $A$ và song song với $(P)$ là $2 x-3 y+z-8=0$}
	 \loigiai{
		\begin{itemchoice}
			\itemch Một vectơ pháp tuyến của mặt phẳng $(P)$  là $\overrightarrow{n}=(2;-3; 1)$.
			\itemch Ta có $d(A,(P))=\dfrac{|2 \cdot 1-3 \cdot(-2)+0-1|}{\sqrt{2^2+(-3)^2+1^2}}=\dfrac{\sqrt{14}}{2}$. 
			\itemch Thay tọa độ $B(1; 1; 1)$ vào phương trình $(P)$  ta được $2\cdot 1-3\cdot 1+1-1=-1 \neq 0 \Rightarrow B(1; 1; 1) \notin(P)$.
			 \itemch Mặt phẳng song song với $(P)$ có phương trình là $2 x-3 y+z+c=0,(c \neq-1)$.\\
			Vì mặt phẳng  đi qua điểm $A(1;-2; 0)$ nên $2\cdot 1-3\cdot (-2)+0+c=0 \Rightarrow c=-8$.\\
			Vậy mặt phẳng cần tìm $2 x-3 y+z-8=0$. 
		\end{itemchoice}
	 }
\end{ex}

% nguồn: Phan1_Lop_12_Hinh_hoc_Chuong_5
\begin{ex}%[2H5V1-3]%[ma:57a9dbf5]
	Trong không gian với hệ tọa độ $Oxyz$,  cho bốn điểm $A(1;2;1)$, $B(2;-1;3)$, $C(2;-1;1)$ và $D(0;3;1)$.
	\choiceTF 
		{\True  Phương trình mặt phẳng $(ABC)$ là $3x+y+5z-10=0$}
		{\True Bốn điểm $A$, $B$, $C$, $D$ tạo thành tứ diện} 
		{Mặt  phẳng $(P)$ chứa $AB$ và song song với $CD$ có một vectơ pháp tuyến là \break $\overrightarrow{a}=(4;-2;7)$}
		{\True Có $2$ mặt phẳng đi qua $2$  điểm $A$, $B$ sao cho khoảng cách từ $C$ và $D$ đến mặt phẳng đó bằng nhau}
	\loigiai{ 
		\begin{itemchoice}
			\itemch Ta có $\overrightarrow{AB} = (1;-3;2)$, $\overrightarrow{AC} = (1;-3;0)$  suy ra $[\overrightarrow{AB}, \overrightarrow{AC}] = (6;2;10)$.\\
			Suy ra  phương trình mặt phẳng $(ABC)$ là
			$$6(x-1)+2(y-2)+10(z-1)=0  \Leftrightarrow 3x+y+5z-10=0.$$
			\itemch Ta có $3(0) + 3 +  5(1) - 10 = -2 \ne 0$ nên $D \notin (ABC)$.\\
			Do  đó bốn điểm $A$, $B$, $C$, $D$ tạo thành tứ diện.
			 \itemch Mặt phẳng $(P)$ chứa $AB$ và song song với $CD$ có một vectơ pháp tuyến là $\overrightarrow{n} = [\overrightarrow{AB}, \overrightarrow{CD}]$.\\
			$\overrightarrow{AB}  = (1;-3;2)$, $\overrightarrow{CD} = (-2;4;0)$.\\
			Suy ra $[\overrightarrow{AB}, \overrightarrow{CD}] =  (-8;-4;-2)$.\\
			Vì $\overrightarrow{n}$ không cùng phương với $\overrightarrow{a}=(4;-2;7)$ nên $\overrightarrow{a}$  không là vectơ pháp tuyến của $(P)$.
			\itemch Trường hợp 1: $(P)$ qua hai điểm $A$, $B$ và song song với $CD$, khi đó $(P)$ có vectơ pháp tuyến là $[\overrightarrow{AB},  \overrightarrow{CD}] = (-8;-4;-2)$ (hoặc $(4;2;1)$).\\
			Phương trình mặt  phẳng $(P)$ là $4(x-1)+2(y-2)+1(z-1)=0 \Leftrightarrow 4x+2y+z-7=0$.\\
			Khoảng cách  từ $C(2;-1;1)$ đến $(P)$ là $\mathrm{d}(C,(P)) = \dfrac{|4\cdot 2+2\cdot
			(-1)+1\cdot 1-7|}{\sqrt{4^2+2^2+1^2}}=0$.\\ 
			Khoảng  cách từ $D(0;3;1)$ đến $(P)$ là $\mathrm{d}(D,(P)) = \dfrac{|4\cdot 0+2\cdot 3+1\cdot 1-7|}{\sqrt{4^2+2^2+1^2}}= 0$.\\
			 Vậy $\mathrm{d}(C,(P)) = \mathrm{d}(D,(P)) = 0$.\\
			Trường hợp 2: $(Q)$ qua hai điểm $A$, $B$ cắt $CD$ tại  trung điểm $I$ của đoạn $CD$.\\
			Trung điểm $I$ của $CD$ là $I(1;1;1)$.\\ 
			Ta có $\overrightarrow{AB} = (1;-3;2)$, $\overrightarrow{AI} = (0;-1;0)$.\\ 
			Vectơ pháp tuyến của $(Q)$ là  $[\overrightarrow{AB}, \overrightarrow{AI}] = (2;0;1)$.\\
			Phương  trình mặt phẳng $(Q)$ là $2(x-1)+0(y-1)+1(z-1)=0 \Leftrightarrow 2x+z-3=0$.\\
			Khoảng cách  từ $C(2;-1;1)$ đến $(Q)$ là $\mathrm{d}(C,(Q)) = \dfrac{|2\cdot 2+1-3|}{\sqrt{2^2+0^2+1^2}}= \dfrac{2}{\sqrt{5}}$.\\
			Khoảng cách từ $D(0;3;1)$ đến $(Q)$ là $\mathrm{d}(D,(Q)) = \dfrac{|2\cdot  0+1-3|}{\sqrt{2^2+0^2+1^2}}= \dfrac{2}{\sqrt{5}}$.\\
			Vậy  $\mathrm{d}(C,(Q)) = \mathrm{d}(D,(Q))$.
		\end{itemchoice}
	 }
\end{ex}

% nguồn: Phan1_Lop_12_Hinh_hoc_Chuong_5
\begin{ex}%[2H5V1-3]%[ma:a5e3b1f8]
	Trong không gian với hệ tọa độ $Oxyz$, cho điểm  $A(1;2;5)$ và mặt phẳng $(\alpha)\colon x + 2y + 2z - 6 = 0$.
	 \choiceTF
		{\True Vectơ $\overrightarrow{n} =  (1;2;2)$ là một vectơ pháp tuyến của $(\alpha)$}
		{Phương trình  mặt phẳng $(\beta)$ đi qua điểm $A$ và song song với mặt phẳng $(\alpha)$ có phương trình $x + 2y + 2z + 15 = 0$}
		{\True Phương trình mặt phẳng  $(\gamma)$ đi qua hai điểm $O$ và $A$ đồng thời vuông góc với mặt phẳng $(\alpha)$ có phương trình $2x - y = 0$}
		{\True Điểm $M \in (\alpha)$ sao cho $A, O, M$  thẳng hàng thì tọa độ $M\left(\dfrac{2}{5};\dfrac{4}{5};2\right)$}
	 \loigiai{
		\begin{itemchoice}
			\itemch Theo định nghĩa vectơ  $\overrightarrow{n} = (1;2;2)$ là một vectơ pháp tuyến của $(\alpha)$.
			\itemch Phương trình mặt phẳng $(\beta)$ đi qua $A$ và song song với $(\alpha)$ nên 
			$\overrightarrow{n_{(\beta)}} = \overrightarrow{n_{(\alpha)}} = (1;2;2)$.\\ 
			Do đó $(\beta)$ có phương trình $x - 1 + 2(y - 2) + 2(z  - 5) = 0 \Leftrightarrow x + 2y + 2z - 15 = 0$.
			\itemch Phương trình  mặt phẳng $(\gamma)$ đi qua hai điểm $O$ và $A$ đồng thời vuông góc với mặt phẳng $(\alpha)$ nên
			$\overrightarrow{OA} =  (1;2;5)$
			$\overrightarrow{n_{(\alpha)}} =  (1;2;2)$
			$\overrightarrow{n_{(\gamma)}} = [\overrightarrow{OA},  \overrightarrow{n_{(\alpha)}}] = (-6;3;0) = -3(2;-1;0)$.\\
			Do đó $(\gamma)$ có phương  trình $2x - y = 0$.
			\itemch Điểm $M \in (\alpha)$ sao cho $A$, $O$, $M$  thẳng hàng.\\
			 Do $M \in (\alpha)\colon x - 6 - 2y - 2z \Rightarrow M(6 - 2b - 2c; b; c)$.\\
			Vì $A,  O, M$ thẳng hàng nên $\overrightarrow{OM} = (6 - 2b - 2c; b; c) = k\overrightarrow{OA} = (k; 2k; 5k)$.
			$$ \heva{&6 - 2b - 2c = k \\& b = 2k \\&c = 5k} \Rightarrow \heva{&6 - 2(2k) - 2(5k) = k \\&b = 2k  \\&c = 5k} \Rightarrow \heva{&6 = 15k \\&b = 2k \\&c = 5k } \Rightarrow \heva{&k = \frac{6}{15} = \dfrac{2}{5} \\&b = 2 \cdot \dfrac{2}{5} = \dfrac{4}{5} \\&c = 5 \cdot \dfrac{2}{5} = 2.} $$
			 Vậy $M\left(\dfrac{2}{5};\dfrac{4}{5};2\right)$.
		\end{itemchoice}
	} 
\end{ex}

% nguồn: Phan1_Lop_12_Hinh_hoc_Chuong_5
\begin{ex}%[2H5V1-3]%[ma:36ce3c0d]
	Trong không gian $Oxyz$, cho  ba điểm $A(1;2;-1)$, $B(2;-1;3)$, $C(-4;7;5)$.
	 \choiceTF
		{\True Mặt phẳng $(ABC)$ có  phương trình $19x+13y+5z-40=0$}
		{Điểm $A(1;2;-1)$ thuộc mặt  phẳng $2x-3y+5z-9=0$}
		{Cho điểm $M$ thỏa mãn $2MA=3MB$. Giá trị lớn nhất của $MC$ bằng $16{,}8$ (làm tròn kết quả đến hàng phần chục)} 
		{\True $G\left(-\dfrac{1}{3}; \dfrac{8}{3}; \dfrac{7}{3}\right)$ là trọng tâm của tam giác $ABC$} 
	\loigiai{ 
		\begin{itemchoice}
			\itemch  Ta có $\overrightarrow{AB}=(1;-3;4)$; $\overrightarrow{AC}=(-6;8;2)$. \\
			Vì  $\overrightarrow{AB}$, $\overrightarrow{AC}$ không cùng phương nên chúng là cặp VTCP của $(ABC)$. \\
			Ta có  $\left[\overrightarrow{AB}, \overrightarrow{AC}\right]=(-38;-26;-10)=-2(19;13;5)$. \\
			Mặt phẳng $(ABC)$ đi qua $A(1;2;-1)$ và nhận $\overrightarrow{n}=(19;13;5)$ làm VTPT, có phương trình  là
			\begin{eqnarray*}
				 &&19(x-1)+13(y-2)+5(z+1)=0 \\
				&\Leftrightarrow& 19x-19+13y-26+5z+5=0  \\
				&\Leftrightarrow&  19x+13y+5z-40=0.
			\end{eqnarray*}
			\itemch Thay $x=1$, $y=2$, $z=-1$ vào phương trình  $2x-3y+5z-9=0$, có
			\begin{eqnarray*}
				&&2 \cdot 1 -3  \cdot 2 + 5 \cdot (-1)-9=0 \\
				&\Leftrightarrow& -18 = 0 \text{ (vô lí)}. 
			\end{eqnarray*}
			Vậy $A(1;2;-1)$ không  thuộc mặt phẳng $2x-3y+5z-9=0$.
			\itemch  Gọi $M(x;y;z)$.\\
			Ta có 
			\begin{eqnarray*}
				 &&2MA=3MB \\
				&\Leftrightarrow& 4\left[(x-1)^2+(y-2)^2+(z+1)^2\right]=9\left[(x-2)^2+(y+1)^2+(z-3)^2\right]  \\
				&\Leftrightarrow&  x^2+y^2+z^2-\dfrac{28}{5}x+\dfrac{34}{5}y-\dfrac{62}{5}z+\dfrac{102}{5}=0.
			\end{eqnarray*}
			Ta có $a=\dfrac{14}{5}$, $b=\dfrac{17}{5}$,  $c=\dfrac{31}{5}$, $d=\dfrac{102}{5}$.\\
			Vì $\left(\dfrac{14}{5}\right)^2 + \left(\dfrac{17}{5}\right)^2  + \left(\dfrac{31}{5}\right)^2 - \dfrac{102}{5} = \dfrac{936}{25} > 0$ nên phương trình trên là phương trình mặt cầu $(S)$ tâm $I\left(\dfrac{14}{5}; \dfrac{17}{5}; \dfrac{31}{5}\right)$, bán kính $R=\dfrac{6\sqrt{26}}{5}$.\\
			Suy ra $M \in  (S)$. \\
			Giá trị lớn nhất của $MC=CI+R=\sqrt{\left(\dfrac{17}{5}+4\right)^2 + \left(\dfrac{17}{5}-7\right)^2 + \left(\dfrac{31}{5}-5\right)^2}  + \dfrac{6\sqrt{26}}{5} \approx 13{,}9$.
			\itemch Đúng.  Tọa độ trọng tâm của tam giác $ABC$ là
			\[G=\left(\dfrac{1+2+(-4)}{3}; \dfrac{2+(-1)+7}{3}; \dfrac{(-1)+3+5}{3}\right)=\left(-\dfrac{1}{3}; \dfrac{8}{3};  \dfrac{7}{3}\right).\]
		\end{itemchoice}
	 }
\end{ex}

% nguồn: Phan1_Lop_12_Hinh_hoc_Chuong_5
\begin{ex}%[2H5V1-3]%[ma:fa1bc5a1]
	Trong không gian $Oxyz$, cho điểm $M(2;3;-1)$ và mặt phẳng $(P)\colon x+2y-3z+1=0$. 
	\choiceTF 
		{\True Véctơ pháp tuyến  của mặt phẳng $(P)\colon x+2y-3z+1=0$ là $\overrightarrow{n}=(1;2;-3)$}
		{Phương trình mặt phẳng $(Q)$ qua điểm $M(2;3;-1)$ và song song  với mặt phẳng $(P)$ là $x+2y-3z+11=0$}
		{\True Phương trình mặt phẳng $(R)$ đi qua $M(2;3;-1)$, song song với trục $Oy$ và vuông góc với  mặt phẳng $(P)$ là $3x+z-5=0$}
		{Cho mặt phẳng $\alpha$ đi qua $M(2;3;-1)$ cắt trục $Ox$, $Oy$, $Oz$ lần lượt tại $3$ điểm $A$, $B$, $C$ sao cho điểm $M$ là trực  tâm của tam giác $ABC$. Phương trình mặt phẳng $(\alpha)$ là $2x+3y+z+14=0$}
	 \loigiai{
		\begin{itemchoice}
			\itemch Véctơ pháp tuyến  của mặt phẳng $(P)$ là $\overrightarrow{n}=(1;2;-3)$.
			\itemch Phương trình mặt phẳng $(Q)$ qua điểm $M(2;3;-1)$ và song song với  mặt phẳng $(P)$ là $x+2y-3z-11=0$.
			\itemch Mặt  phẳng $(R)$ song song $Oy$ và vuông góc với mặt phẳng $(P)$ nên có cặp véctơ chỉ phương là $\overrightarrow{j}=(0;1;0)$ và $\overrightarrow{n}=(1;2;-3)\Rightarrow \overrightarrow{n_{(R)}}=\left[\overrightarrow{j},\overrightarrow{n}\right]=(-3;0;-1)$.\\
			Mặt phẳng $(R)$ đi qua điểm $M(2;3;-1)$ và có một véctơ pháp tuyến  $\overrightarrow{n}_{(R)}=(-3;0;-1)$ có phương trình là $3x+z-5=0$.
			\itemch Do mặt phẳng $\alpha$ đi qua $M(2;3;-1)$ cắt trục $Ox$, $Oy$, $Oz$ lần lượt tại $3$ điểm $A$, $B$, $C$ sao cho điểm $M$ là trực tâm của tam giác $ABC$ nên  $\overrightarrow{OM}=(2;3;-1)$ là véctơ pháp tuyến của mặt phẳng $(\alpha)$. Vậy phương trình mặt phẳng $(\alpha)$ là $2x+3y-1-z-14=0$.
		\end{itemchoice}
	} 
\end{ex}

% nguồn: Phan1_Lop_12_Hinh_hoc_Chuong_5
\begin{ex}%[2H5V1-3]%[ma:e79c55d0]
	Trong không gian với  hệ trục tọa độ $Oxyz$, cho ba điểm $A(1;2;-2)$; $B(2;1;2)$; $C(3;-2;1)$.
	 \choiceTF
		{\True Mặt phẳng $(P)$ đi qua ba điểm $A$, $B$, $C$ có một  vectơ pháp tuyến là $\overrightarrow{n} = \left[\overrightarrow{AB}, \overrightarrow{AC}\right]$}
		{Phương trình mặt phẳng $(P)$ là $13x + 5y - 2z + 27  = 0$}
		{Phương trình mặt phẳng đi qua điểm $A(1;2;-2)$ và nhận  $\overrightarrow{BC}$ làm vectơ pháp tuyến có dạng $1(x + 1) - 3(y - 2) - 1(z - 2) = 0$}
		{\True Phương trình mặt phẳng đi qua ba điểm $E;F;K$ lần lượt là hình chiếu vuông góc của $A$ lên trục $Ox$; $Oy$; $Oz$ có  phương trình $\dfrac{x}{1} + \dfrac{y}{2} + \dfrac{z}{-2} = 1$}
	 \loigiai{
		\begin{itemchoice}
			\itemch Vì $\overrightarrow{AB}=(1;-1;4)$  và $\overrightarrow{AC}=(2;-4;3)$ không cùng phương (do $1/2 \neq -1/-4$).\\
			Mặt phẳng $(P)$ đi qua ba điểm  $A,B,C$ có một vectơ pháp tuyến là $\overrightarrow{n} = [\overrightarrow{AB}, \overrightarrow{AC}]$.
			\itemch Ta có  $\overrightarrow{AB} = (1;-1;4)$, $\overrightarrow{AC} = (2;-4;3)$.\\
			$\overrightarrow{n} =\left[\overrightarrow{AB}, \overrightarrow{AC}\right] = ((-1) \cdot 3 - 4 \cdot (-4); 4 \cdot 2 - 1 \cdot 3; 1 \cdot (-4) - (-1) \cdot  2) = (13; 5; -2)$.\\
			Phương trình mặt phẳng $(P)$ đi qua $A(1;2;-2)$ và nhận $\overrightarrow{n}=(13;5;-2)$  làm vectơ pháp tuyến là
			$13(x-1) + 5(y-2) - 2(z+2) = 0\Leftrightarrow 13x + 5y -  2z - 27 = 0$.
			\itemch Phương trình  mặt phẳng đi qua $A(1;2;-2)$ và nhận $\overrightarrow{BC}= (1;-3;-1)$ làm vectơ pháp tuyến là $1(x-1) - 3(y-2) - 1(z+2) = 0$.
			\itemch Tọa độ ba điểm $E(1;0;0)$, $F(0;2;0)$, $K(0;0;-2)$ lần lượt là  hình chiếu vuông góc của $A(1;2;-2)$ lên trục $Ox$, $Oy$, $Oz$.\\
			Mặt phẳng $(EFK)$ có phương trình theo đoạn chắn là $\dfrac{x}{1} + \dfrac{y}{2} +  \dfrac{z}{-2} = 1$.
		\end{itemchoice}
	 }
\end{ex}

% nguồn: Phan1_Lop_12_Hinh_hoc_Chuong_5
\begin{ex}%[2H5V1-3]%[ma:b3545750]
	Trong không gian $Oxyz$, cho bốn điểm $A(2;-1;0)$, $B(1;1;3)$, $C(-3;2;1)$,  $D(-1;14)$.
	 \choiceTF
		{\True Mặt phẳng $(P)$  qua $A,B$ song song $CD$ có một vectơ pháp tuyến là $\overrightarrow{n}=(3;3;-1)$}
		{Mặt phẳng $(P)$ đi qua  điểm $M(2;-1;5)$}
		{\True Gọi $(Q)$ là mặt phẳng đi qua $C,D$ song song với $AB$.  Khi đó $\mathrm{d}(A,(Q)) = \dfrac{7}{\sqrt{19}}$}
		{\True Mặt phẳng cách  đều $AB,CD$ có dạng $(R)\colon 3x+3y-z+ \dfrac{1}{2}=0$}
	 \loigiai{
		\begin{itemchoice}
			\itemch \textbf{Đúng}. Mặt phẳng $(P)$ qua $A$, $B$ song song $CD$ nên mặt phẳng $(P)$ nhận hai vectơ $\overrightarrow{AB}=(-1;2;3)$, $\overrightarrow{CD}=(2;-1;3)$ làm cặp vectơ chỉ phương suy ra mặt phẳng $(P)$ có vectơ  pháp tuyến $\overrightarrow{n}=\left[\overrightarrow{AB},\overrightarrow{CD}\right]=(9;9;-3)$.
			\itemch \textbf{Sai}. Phương trình mặt phẳng $(P)\colon 3x+3y-z-3=0$ thay toạ độ điểm $M$  vào phương trình mặt phẳng $(P)$ ta thấy không thỏa mãn.
			 \itemch \textbf{Đúng}. $\mathrm{d}(A,(Q))=\mathrm{d}(C,(P))= \dfrac{7}{\sqrt{19}}$.
			\itemch \textbf{Đúng}. Mặt phẳng cách  đều $AB,CD$ song song với $(P)$ nên có phương trình dạng $(R)\colon 3x+3y-z+d=0$.\\
			Gọi $H$ là trung điểm của $AC \Rightarrow  H\left( \dfrac{-1}{2}; \dfrac{1}{2}; \dfrac{1}{2}\right)$, suy ra $H \in (R) \Rightarrow d = \dfrac{1}{2}$.\\
			Vậy $\left(  R \right)\colon 3x+3y-z+ \dfrac{1}{2}=0$.
		\end{itemchoice}
	 }
\end{ex}

% nguồn: Phan1_Lop_12_Hinh_hoc_Chuong_5
\begin{ex}%[2H5V1-3]%[ma:5da70d9e]
	Trong không gian $Oxyz$, cho mặt cầu $(S)$ có  tâm $I(2; 3;-1)$ cắt đường thẳng $(d)\colon \dfrac{x-11}{2}=\dfrac{y}{1}=\dfrac{z+25}{-2}$ tại hai điểm $A$, $B$ sao cho $AB=16$. Xét tính đúng sai của các khẳng định sau
	\choiceTF 
		{\True Đường thẳng $(d)$ có  vectơ chỉ phương là $\vec{u}=(2; 1;-2)$}
		{Đường thẳng $(d)$  đi qua điểm $A(5;-3;-31)$}
		{\True Mặt phẳng $(P)$ chứa $I(2; 3;-1)$  và vuông góc với đường thẳng $(d)$ có phương trình là $2x+y-2z-9=0$}
		{Mặt cầu $(S)$ có phương  trình là $(x-2)^2+(y-3)^2+(z+1)^2=225$}
	 \loigiai{
		\begin{itemchoice}
			\itemch Đường thẳng ($d$) có vectơ chỉ phương  là $\vec{u}=(2; 1;-2)$.
			\itemch Thay $A(5;-3;-31)$  vào phương trình đường thẳng $(d)$ ta có $$\dfrac{5-11}{2}=\dfrac{-3}{1} \neq \dfrac{-31+25}{-2}$$
			nên  điểm $A$ không thuộc đường thẳng $(d)$.
			\itemch Vì mặt phẳng $(P)$ chứa $I(2; 3;-1)$ và vuông góc với đường thẳng $(d)$ nên $(P)$ nhận vectơ chỉ  phương của $(d)$ là $\vec{u}=(2; 1;-2)$ làm vectơ pháp tuyến.\\
			Khi đó, phương trình  mặt phẳng $(P)$ là $2(x-2)+1(y-3)-2(z+1)=0$ hay $2x+y-2z-9=0$.
			\itemch Gọi $H$ là hình chiếu của $I$ lên đường thẳng ($d$) nên $IH$ và ($d$) vuông góc với nhau.\\ 
			Gọi mặt phẳng $(P)$ chứa $I(2; 3;-1)$ và vuông góc với đường thẳng  $(d)$ có phương trình là $2x+y-2z-9=0$.\\
			Suy ra, $H$ là giao điểm của đường thẳng $(d)$ và mặt  phẳng $(P)$.\\
			 Vì $H$ thuộc đường thẳng ($d$) nên $H(2t+11, t,-2t-25)$.\\
			Mặt khác, $H$ thuộc mặt phẳng $(P)$ nên tọa độ $H(2t+11, t,-2t-25)$ thỏa mãn  $2(2t+11)+t-2(-2t-25)-9=0$ hay $t=-7$.\\
			Khi đó, điểm $H(-3,-7,-11)$ và  $IH=\sqrt{(-3-2)^2+(-7-3)^2+(-11+1)^2}=15$
			Ta được $R=\sqrt{IH^2+\left(\dfrac{AB}{2}\right)^2}=\sqrt{15^2+\left(\dfrac{16}{2}\right)^2}=17$.\\ 
			 Vậy phương trình mặt cầu $(S)$ có tâm $I(2; 3;-1)$ và bán kính $R=17$ là
			$$(x-2)^2+(y-3)^2+(z+1)^2=17^2=289.$$ 
		\end{itemchoice}
	 }
\end{ex}

% nguồn: Phan1_Lop_12_Hinh_hoc_Chuong_5
\begin{ex}%[2H5V1-3]%[ma:cd4f06a6]
	Trong không gian với hệ tọa độ $Oxyz$, cho $A(0; 1; 1), B(1; 0;-3), C(-1;-2;-3)$ và mặt cầu $(S)$ có phương trình  $x^2+y^2+z^2-2x+2z-2=0$. Xét tính đúng sai của các khẳng định sau
	 \choiceTF
		{\True Mặt cầu $(S)$ có bán kính  $R=2$}
		{Mặt phẳng  $(ABC)$ có phương trình $2x-2y+z-1=0$}
		{\True Mặt  phẳng $(ABC)$ cắt mặt cầu $(S)$ theo một đường tròn có bán kính bằng $\dfrac{4\sqrt{2}}{3}$}
		{\True Điểm  $D\ (a; b; c)$ thuộc mặt cầu $(S)$ sao cho thể tích tứ diện $ABCD$ lớn nhất. Khi đó, $a+b+c=\dfrac{2}{3}$}
	 \loigiai{
		\begin{itemchoice}
			\itemch Mặt cầu $(S)$ có  tâm $I(1; 0;-1)$, bán kính $R=2$.
			\itemch $\overrightarrow{AB}=(1;-1;-4), \overrightarrow{AC}=(-1;-3;-4)$ nên mặt phẳng ($ABC$) có một vectơ  pháp tuyến là $\vec{n}=\left[\overrightarrow{AB}, \overrightarrow{AC}\right]=(-8; 8;-4)$.\\
			Suy ra mặt phẳng  ($ABC$) có phương trình $2x-2y+z+1=0$.
			\itemch Ta có  $\mathrm{d}(I,(ABC.)=\dfrac{|2\cdot 1-2\cdot 0-1+1|}{\sqrt{2^2+(-2)^2+1}}=\dfrac{2}{3}$.\\
			Vậy bán kính đường tròn giao tuyến của mặt phẳng $(ABC)$ và mặt cầu $(S)$ là  $r=\sqrt{R^2-\left[\mathrm{d}(I,(ABC))\right]^2}=\sqrt{2^2-\left(\dfrac{2}{3}\right)^2}=\dfrac{4\sqrt{2}}{3}$.
			\itemch Ta có $V_{ABCD}=\dfrac{1}{3} \mathrm{d}(D,(ABC)) \cdot S_{ABC}$ nên $V_{ABCD}$ lớn nhất khi và chỉ  khi $\mathrm{d}(D,(ABC))$ lớn nhất.\\ Gọi $\Delta$ là đường thẳng qua điểm $I$ và vuông góc với mặt phẳng ($ABC$).\\
			Suy ra $\Delta\colon\heva{&x=1+2t \\&y=-2t  \\ &z=-1+t}$.\\
			Gọi $D_1; D_2$ là các giao điểm của  $\Delta$ và mặt cầu $(S)$. Tọa độ điểm $D_1; D_2$ thỏa mãn hệ phương trình: \allowdisplaybreaks
			\begin{eqnarray*}
				&&\heva{&x=1+2t \\&y=-2t \\&z=-1+t  \\&x^2+y^2+z^2-2x+2z-2=0}\\ &\Rightarrow&\hoac{&t=\dfrac{2}{3} \\&t=\dfrac{-2}{3}}\\
				&\Rightarrow& D_1\left(\dfrac{7}{3}; \dfrac{-4}{3}; \dfrac{-1}{3}\right); D_2\left(\dfrac{-1}{3};  \dfrac{4}{3}; \dfrac{-5}{3}\right)
			\end{eqnarray*}
			Ta  thấy: $\mathrm{d}\left(D_1,(ABC)\right) > d\left(D_2,(ABC)\right)$.\\ Vậy điểm $D\left(\dfrac{7}{3};-\dfrac{4}{3};-\dfrac{1}{3}\right) \Rightarrow a+b+c=\dfrac{2}{3}$.
		\end{itemchoice}
	} 
\end{ex}


%%%%%%%%%%%%%%%% TRẢ LỜI NGẮN %%%%%%%%%%%%%%%%

% nguồn: Phan2_Lop_12_Hinh_hoc_Chuong_5
\begin{bt}%[2H5N1-3]%[ma:46e779f1]
	Trong không gian với hệ tọa độ $Oxy$, cho điểm $M(2;1;-1)$ và mặt phẳng $(P) \colon ax-3y-z-8=0$ (với $a$ là hằng số). Biết rằng  mặt phẳng $(P)$ đi qua điểm $M$. Giá trị của $a$ bằng bao nhiêu?
	\shortans{5} 
	\loigiai{Vì $M(2;1;-1) \in  (P) \colon ax-3y-z-8=0$
		nên $$a \cdot 2-3 \cdot 1 -(-1)-8=0 \Leftrightarrow 2a-10=0  \Leftrightarrow a=5.$$
	 }
\end{bt}

% nguồn: Phan2_Lop_12_Hinh_hoc_Chuong_5
\begin{bt}%[2H5H1-3]%[ma:51a93bb8]
	Trong không gian tọa độ $Oxyz$, phương trình mặt phẳng $(P)$ đi qua ba điểm $A(1; 0; -2)$, $B(-3; 1; 1)$, $C(5; 5; -5)$ có dạng $ax + by + cz + 5 = 0$, khi đó giá trị của biểu thức $a + b + c$ bằng 
	 \shortans{7}
	 \loigiai{
		Ta có $\overrightarrow{AB} = (-4; 1; 3), \ \overrightarrow{AC} = (4; 5; -3) \Rightarrow [\overrightarrow{AB}, \overrightarrow{AC}] = (-18; 0;  -24)$.\\
		Một  vectơ pháp tuyến của mặt phẳng là $\overrightarrow{n} = -\dfrac{1}{6}[\overrightarrow{AB}, \overrightarrow{AC}] = (3; 0; 4)$.\\
		Phương trình mặt phẳng  $(P)$ đi qua điểm $A(1; 0; -2)$ là
		$$3(x - 1) + 4(z  + 2) = 0 \Rightarrow 3x + 4z + 5 = 0.$$
		Suy ra $a = 3$, $b = 0$, $c = 4 \Rightarrow a + b +  c = 7$.
	 }
\end{bt}

% nguồn: Phan2_Lop_12_Hinh_hoc_Chuong_5
\begin{bt}%[2H5H1-3]%[ma:c46c737f]
	Trong không gian với hệ tọa độ $(Oxyz)$, mặt phẳng $x+my+nz=0$ ($m$, $n$ là các số thực) đi qua hai điểm $A(2;3;1)$ và $B(4;1;7)$. Giá trị $5m-6n$ bằng  bao nhiêu?
	\shortans{0,5} 
	\loigiai{Mặt phẳng $x+my+nz=0$ ($m$, $n$ là các số thực) đi  qua hai điểm $A(2;3;1)$ và $B(4;1;7)$ nên ta có
		\[\heva{&2+3m+n=0\\&4+m+7n=0} \Leftrightarrow \heva{&3m+n=-2\\&m+7n=-4} \Leftrightarrow \heva{&m=\dfrac{-1}{2}\\&n=\dfrac{-1}{2}.}\] 
		Do đó $5m-6n=5  \cdot \left(\dfrac{-1}{2}\right)-6 \cdot \left(\dfrac{-1}{2}\right)=\dfrac{1}{2}=0{,}5$.
	} 
\end{bt}

% nguồn: Phan2_Lop_12_Hinh_hoc_Chuong_5
\begin{bt}%[2H5H1-3]%[ma:f33ae897]
	Trong một khung lưới ô vuông gồm các hình lập phương, người ta đưa ra một cách kiểm tra bốn nút lưới (\textit{đỉnh hình lập phương}) bất kì có đồng phẳng  hay không bằng cách gắn hệ trục tọa độ $Oxyz$ vào khung lưới ô vuông và lập phương trình mặt phẳng đi qua ba nút lưới trong bốn nút lưới đã cho.
	Giả sử có ba nút lưới mà tọa độ  lần lượt là $(1;1;10)$, $(4;3;1)$, $(3;2;5)$ và mặt phẳng đi qua ba nút lưới đó có phương trình $x+my+nz+p=0$.
	Giá trị của  $m+n+p$ là bao nhiêu?
	 \shortans{-10}
	\loigiai{ 
		Xét các điểm $A(1;1;10)$, $B(4;3;1)$, $C(3;2;5)$. 
		Khi đó, ta  có
		 \[
		 \overrightarrow{AB}=(3;2;-9),\quad \overrightarrow{AC}=(2;1;-5)
		 \Rightarrow [\overrightarrow{AB},\overrightarrow{AC}]=(-1;-3;-1).
		 \]
		Gọi $\vec{n}$ là một  vectơ pháp tuyến của mặt phẳng đi qua ba điểm $A,B,C$.
		Chọn $\vec{n}=(1;3;1)$. 
		Suy  ra phương trình mặt phẳng đi qua ba nút lưới là:
		 \[
		 1(x-1)+3(y-1)+1(z-10)=0
		 \Rightarrow x+3y+z-14=0.
		\] 
		 Vậy $m=3$, $n=1$, $p=-14$ nên $m+n+p=-10$.
	} 
\end{bt}

% nguồn: Phan2_Lop_12_Hinh_hoc_Chuong_5
\begin{bt}%[2H5H1-3]%[ma:e01f2a33]
	Trong không gian $Oxyz$, cho hai điểm $A(-1;2;0)$ và $B(3;0;2)$. Mặt phẳng trung trực của đoạn thẳng $AB$ có dạng $ax+by+cz-2=0$. Tính  giá trị biểu thức $K=a-b+c$.
	\shortans{4} 
	 \loigiai{
		Gọi  $I$ là trung điểm của đoạn thẳng $AB$, khi đó tọa độ của điểm $I$ là
		\begin{center}
			 $\heva{&x_I=\dfrac{x_A+x_B}{2}=\dfrac{-1+3}{2}=1\\&y_I=\dfrac{y_A+y_B}{2}=\dfrac{2+0}{2}=1\\&z_I=\dfrac{z_A+z_B}{2}=\dfrac{0+2}{2}=1.}$
		\end{center}
		 Suy ra $I(1;1;1)$.\\
		Ta có $\overrightarrow{AB}=(4;-2;2)$.\\ 
		Gọi $(P)$ là mặt phẳng trung trực của đoạn thẳng $AB$. Suy ra $(P)$ đi qua điểm $I$ và nhận véc-tơ  $\overrightarrow{AB}$ là một véc-tơ pháp tuyến. Phương trình của mặt phẳng $(P)$ là
		\begin{center}
			$4(x-1)-2(y-1)+2(z-1)=0 \Leftrightarrow 4x-2y+2z-4=0 \Leftrightarrow 2x-y+z-2=0$. 
		\end{center}
		 Suy ra $a=2$, $b=-1$, $c=1$.\\
		Do đó $K=a-b+c=2-(-1)+1=4$. 
	} 
\end{bt}

% nguồn: Phan2_Lop_12_Hinh_hoc_Chuong_5
\begin{bt}%[2H5H1-3]%[ma:0424e3df]
	\immini{ 
	Cho hình minh họa một nhà kho trong không gian với hệ tọa độ $Oxyz$ (đơn vị trên mỗi trục là mét). Hai mái $EFIK$, $HGIK$ có kích thước bằng nhau. Biết rằng chiều cao của nhà kho là $9$ m và các bức tường của nhà kho tạo thành hình hộp chữ nhật $ABCD.EFGH$ với $AB=10$ m, $AD=24$ m, $AE=7$ m. Biết mặt phẳng $(EFIK)$ có phương trình $ax+by+cz+42=0$. Tính giá trị của biểu thức $a-b+c$. 
	 }{
	\begin{tikzpicture}[scale=0.8,font=\footnotesize,line join=round,line cap=round,>=stealth]
			 \path
				(0,0) coordinate  (O) node[left]{$B$}
				(4,0) coordinate (A) node[below]{$C$} 
				(0,3) coordinate (E)  node[left]{$F$}
				($(A)+(E)-(O)$) coordinate (F) node[right]{$G$} 
				(A)+(1,2) coordinate (B)  node[below right]{$D$}
				($(O)+(B)-(A)$) coordinate (C) node[above  right]{$O\equiv A$}
				($(C)+(E)-(O)$) coordinate (H) node[above left]{$E$} 
				 ($(H)+(B)-(C)$) coordinate (G) node [right]{$H$}
				($(O)!.6!(A)$) coordinate  (x)
				(x)+(0,4.5)  coordinate (P) node[below]{$I$}
				($(P)+(H)-(E)$) coordinate (Q)  node[above]{$K$};
			\draw[->]  (H)--($(H)+(90:1)$) node[above]{$z$};
			\draw[->]  (O)--($(C)!1.3!(O)$) node[below]{$x$};
			\draw[dashed,->] (C)--($(C)!1.3!(B)$) node[above]{$y$}; 
			\draw (A)--(B)--(G)--(Q)--(H)--(E)--(F)--cycle (E)--(P)--(F) 
				 (P)--(Q)
				(F)--(G) 
				 (A)--(O)--(E);
			\draw[dashed]  (C)--(B)
				 (C)--(H)--(G)
				(C)--(O); 
		\end{tikzpicture}
	} 
	 \shortans{-7}
	\loigiai{ 
		Theo đề bài,  $A(0;0;0)$, $B(10;0;0)$, $D(0;24;0)$, $E(0;0;7)$, $F(10;0;7)$, $G(10;24;7)$, $H(0;24;7)$, $K(5;12;9)$.\\
		Giả sử $(EFIK) \colon  ax+by+cz+42=0$.\\
		Vì $E$, $F$, $K$ là ba điểm thuộc mặt  phẳng $(EFIK)$ nên
		$\heva{&7c + 42 = 0\\&10a + 7c + 42 = 0\\&5a + 12b  + 9c + 42 = 0.}$\\
		Suy ra $a=0$, $b=1$,  $c=-6$.\\
		 Vậy $a-b+c=-7$.
	 }
\end{bt}

% nguồn: Phan2_Lop_12_Hinh_hoc_Chuong_5
\begin{bt}%[2H5H1-3]%[ma:7b8dbc92]
	 Trong không gian $Oxyz$, cho hai điểm $A(2; 0; 1)$ và $2$ vectơ $\overrightarrow{a} = (1; -3; 2)$, $\overrightarrow{b} = (-1; -1; 2)$. Một mặt phẳng $(Q)$ chứa điểm $A$ và nhận $\overrightarrow{a}, \overrightarrow{b}$ làm cặp vecto chỉ phương có phương trình là $ax + by + cz - 3 = 0$. Tính $a + b + c$.
	\shortans{3} 
	 \loigiai{
		Vectơ pháp tuyến của mặt phẳng $(Q)$ là $\overrightarrow{n_{(Q)}} = [\overrightarrow{b},  \overrightarrow{a}] = (4; 4; 4)$.\\
		 Phương trình mặt phẳng $(Q)$ là
		$$4(x -  2) + 4(y - 0) + 4(z - 1) = 0 \Rightarrow 4x + 4y + 4z - 12 = 0 \Rightarrow x + y + z - 3 = 0.$$
		Suy ra $a = b  = c = 1 \Rightarrow a + b + c = 3$.
	} 
\end{bt}

% nguồn: Phan2_Lop_12_Hinh_hoc_Chuong_5
\begin{bt}%[2H5H1-3]%[ma:db1dcc95]
	Trong không gian với hệ tọa độ $Oxyz$, cho  tứ diện $ABCD$ có $A(3;0;0)$, $B(0;6;0)$, $C(0;0;9)$, $D(3;6;9)$. Phương trình mặt phẳng trung trực của đoạn thẳng $CD$ có dạng $ax+by+cz-15=0$ ($a;b;c\in \mathbb{Z}$). Tính giá trị $S=a+b+c$.
	\shortans{6} 
	 \loigiai{
		Ta  có $\vec{CD}=(3;6;0)$.\\
		Gọi $I$ là  tung điểm của $CD$ suy ra $I\left(\dfrac{3}{2};3;9\right)$.\\
		Mặt phẳng trung trực  của đoạn thẳng $CD$ đi qua trung điểm $I$ của $CD$ và vuông góc với $CD$ nên có véc-tơ pháp tuyến $\vec{n}=\dfrac{1}{3}\vec{CD}=(1;2;0)$.\\
		 Phương trình mặt phẳng trung trực của đoạn thẳng $CD$ là
		 $$1\left(x-\dfrac{3}{2}\right)+2(y-3)+0(z-9)=0\Leftrightarrow 2x+4y-15=0.$$
		Khi đó $\heva{&a=2\\&b=4\\&c=0}\Rightarrow S=6$. 
	 }
\end{bt}

% nguồn: Phan2_Lop_12_Hinh_hoc_Chuong_5
\begin{bt}%[2H5V1-3]%[ma:2ead0ca0]
	Trong không gian với hệ trục $Oxyz$, cho mặt phẳng $(\alpha)\colon x + by + cz + d = 0$ vuông góc với mặt phẳng $(\beta)\colon x + 2y + 3z + 4 = 0$ và chứa giao tuyến của hai mặt phẳng $(P)\colon x +  3y + z - 7 = 0$, $(Q)\colon x - y + z + 1 = 0$. Tìm $d$.
	 \shortans{3}
	\loigiai{ 
		Ta có vectơ pháp  tuyến của $(\beta), (P), (Q)$ lần lượt là $\overrightarrow{n}_1=(1;2;3)$, $\overrightarrow{n}_2=(1;3;1)$, $\overrightarrow{n}_3=(1;-1;1)$.\\
		Khi  đó $\overrightarrow{n}_{\alpha} = [\overrightarrow{n}_1, [\overrightarrow{n}_2, \overrightarrow{n}_3]] = (-8;16;-8) = -8(1;-2;1)$.\\
		Gọi  $A(x;y;z) \in d$ là giao tuyến của $(P)$ và $(Q)$, khi đó toạ độ điểm $A$ thoả mãn hệ
		$\heva{& x+3y+z-7=0  \\& x-y+z+1=0}$.\\
		Cho $x=0$ ta có $\heva{&3y+z-7=0 \\&-y+z+1=0} \Leftrightarrow \heva{&y=2  \\&z=1}$ khi đó $A(0;2;1)$.\\
		Do $(\alpha)$ chứa giao tuyến của $(P)$ và $(Q)$ nên  $(\alpha)$ đi qua $A(0;2;1)$.\\
		Phương trình $(\alpha)\colon x - 2(y-2) + z - 1 = 0  \Leftrightarrow x - 2y + z + 3 = 0$. \\
		 Vậy $d=3$.
	} 
\end{bt}

% nguồn: Phan2_Lop_12_Hinh_hoc_Chuong_5
\begin{bt}%[2H5V1-3]%[ma:e1762a93]
	Trong không gian $Oxyz$ cho các điểm $A(2;0;0)$, $B(0;4;0)$, $C(0;0;6)$, $D(2;4;6)$. Gọi $(P)$ là mặt phẳng song song với mặt phẳng $(ABC)$, $(P)$ cách đều $D$ và mặt phẳng $(ABC)$. Phương trình của mặt phẳng $(P)$ là $ax+by+2z+d=0$ với $a$, $b$, $d \in \mathbb{Z}$. Giá trị của  $d$ bằng bao nhiêu?
	 \shortans{-24}
	\loigiai{ 
		Phương trình mặt phẳng $(ABC)$ là $ \dfrac{x}{2} + \dfrac{y}{4} +  \dfrac{z}{6} = 1 \Leftrightarrow 6x+3y+2z-12=0$.\\
		Mặt phẳng $(P)$ song song với mặt phẳng $(ABC)$ nên $(P)$ có dạng $ 6x+3y+2z+d=0 \quad  (d \neq -12)$.\\
		Vì $(P)$ cách đều $D$ và mặt  phẳng $(ABC)$ nên ta có
		\[\mathrm{d}(D;(P)) = \mathrm{d}((ABC);(P)) \Leftrightarrow \mathrm{d}(D;(P)) =  \mathrm{d}(A;(P)) \Leftrightarrow |36+d| = |12+d| \Leftrightarrow d=-24. \]
		Vậy phương trình mặt phẳng  $(P)$ là $6x+3y+2z-24=0$.
	 }
\end{bt}

% nguồn: Phan2_Lop_12_Hinh_hoc_Chuong_5
\begin{bt}%[2H5V1-3]%[ma:e955b9c3]
	Trong không gian $Oxyz$, mặt phẳng $(Q)$ song song với mặt phẳng $(P)\colon 2x  - y - 2z + 19 = 0$ và cách $(P)$ một khoảng bằng $5$ có dạng $ax + y + cz + d = 0$. Tính $a+c+d$.
	\shortans{11} 
	 \loigiai{
		Trên mặt phẳng $(P)\colon 2x - y - 2z + 4 = 0$  chọn điểm $M(0;4;0)$.\\
		Do $(Q)$  song song với mặt phẳng $(P)$ nên $(Q)$ có dạng $2x-y-2z+D=0$ với $D \neq 19$.\\
		Vì $$\mathrm{d}((P),(Q))=5 \Leftrightarrow d(M,(Q)) = 5 \Leftrightarrow  \dfrac{|-4+D|}{\sqrt{2^2+(-1)^2+(-2)^2}} = 5 \Leftrightarrow |-4+D| = 15 \Leftrightarrow
		\hoac{&D=-11 \\&D=19.}$$ 
		Mặt phẳng $(Q)$ thỏa mãn yêu cầu bài toán  $2x-y-2z-11=0$. Vậy $a+c+d = 11$.
	 }
\end{bt}

% nguồn: Phan2_Lop_12_Hinh_hoc_Chuong_5
\begin{bt}%[2H5V1-3]%[ma:0c25c7e1]
	Trong không gian với hệ trục tọa độ $Oxyz$, cho điểm $H(1; 2; 3)$ là trực tâm của tam giác $ABC$ với $A$, $B$, $C$ lần lượt nằm trên các trục $Ox$, $Oy$, $Oz$ (khác gốc tọa độ). Phương trình mặt phẳng đi qua ba điểm $A$, $B$, $C$ có dạng $mx + ny + pz - 14 = 0$ (với $m$, $n$, $p$ $\in \mathbb{Z}$).  Tính $m + n + p$.
	\shortans{6} 
	 \loigiai{
		Gọi  $A(a; 0; 0)$, $B(0; b; 0)$, $C(0; 0; c)$.\\
		Khi đó $\overrightarrow{AH} = (1 - a; 2; 3)$, $\overrightarrow{BH} = (1; 2 - b; 3)$, $\overrightarrow{BC} = (0; -b; c)$, $\overrightarrow{AC} =  (-a; 0; c)$.\\
		 Do $H$ là trực tâm tam giác $ABC$ nên\\
		 $$\heva{
		& \overrightarrow{AH}  \cdot \overrightarrow{BC} = 0 \\
		 & \overrightarrow{BH} \cdot \overrightarrow{AC} = 0
		 }
		 \Rightarrow
		\heva{ 
		&-2b +  3c = 0 \\
		&-a + 3c = 0. 
		 }$$
		Mặt phẳng $(ABC)$ có phương trình $\dfrac{x}{a} + \dfrac{y}{b} + \dfrac{z}{c}  = 1$.\\
		Vì $H \in (ABC)$ nên $\dfrac{1}{a} +  \dfrac{2}{b} + \dfrac{3}{c} = 1$.\\
		Thay $a = 3c$,  $b = \dfrac{3c}{2}$ vào, ta có\\
		$$\dfrac{1}{3c} + \dfrac{2}{\dfrac{3c}{2}} + \dfrac{3}{c} = \dfrac{1}{3c}  + \dfrac{4}{3c} + \dfrac{3}{c} = \dfrac{1 + 4 + 9}{3c} = 1 \Rightarrow c = \dfrac{14}{3}$$
		Suy ra $a =  14$, $b = 7$, $c = \dfrac{14}{3}$.\\
		Phương trình mặt phẳng $\dfrac{x}{14} + \dfrac{y}{7} +  \dfrac{3z}{14} = 1 \Rightarrow x + 2y + 3z - 14 = 0$.\\
		Vậy $m = 1$,  $n = 2$, $p = 3 \Rightarrow m + n + p = 6$.
	 }
\end{bt}

% nguồn: Phan2_Lop_12_Hinh_hoc_Chuong_5
\begin{bt}%[2H5V1-3]%[ma:6855fcc4]
	Một phần sân trường được định vị bởi các điểm $A$, $B$, $C$, $D$  như hình vẽ.
	\begin{center}
		\begin{tikzpicture}[scale=0.8,font=\footnotesize,line join=round,line cap=round,>=stealth]
				\path 
					(0,0) coordinate (A) 
					(9,0) coordinate (B) 
					(9,-6) coordinate (C) 
					(0,-5) coordinate  (D);
				\draw  (A)--(B)--(C)--(D)--cycle;
				\node  [above] at ($(A)!0.5!(B)$) {$25$ m};
				\node [right] at ($(B)!0.5!(C)$)  {$18$ m};
				\node [left] at ($(A)!0.5!(D)$) {$15$ m}; 
				\foreach \i/\g in {A/90,B/90,C/-90,D/-90}{\draw (\i) circle (0pt) 
					($(\i)+(\g:4mm)$)  node[scale=1]{$\i$};}
			\end{tikzpicture}
	\end{center}
	Bước đầu, chúng được lấy \lq\lq thăng bằng\rq\rq\ để  có cùng độ cao. Biết $ABCD$ là hình thang vuông ở $A$ và $B$ với độ dài $AB=25$ m, $AD=15$ m, $BC=18$ m. Do yêu cầu kỹ thuật, khi lát phẳng phần sân trường phải thoát nước về góc sân ở $C$ nên người ta lấy độ cao ở các điểm $B$, $C$, $D$ xuống thấp hơn so với độ cao ở $A$ lần lượt là $10$ cm, $a$ cm, $6$ cm tương ứng sao cho bốn điểm $A$, $B'$, $C'$, $D'$ đồng phẳng. Tính giá trị của $a$.
	 \shortans{17,2}
	 \loigiai{
		Chọn hệ trục toạ độ $Oxyz$ sao cho $O\equiv A$, trục  $Ox\equiv AD$, trục $Oy\equiv AB$.
		 Khi đó:
		 \[
		A(0;0;0),\quad B(0;2500;0),\quad C(1800;2500;0),\quad  D(1500;0;0).
		 \]
		Khi hạ độ cao ở  các điểm $B$, $C$, $D$ xuống thấp hơn so với độ cao ở $A$ lần lượt là $10$ cm, $a$ cm, $6$ cm tương ứng, ta có
		 \[
		B'(0;2500;-10),\quad  C'(1800;2500;-a),\quad D'(1500;0;-6).
		\] 
		Theo giả thiết, bốn điểm $A,B',C',D'$ đồng phẳng.\\  Khi đó mặt phẳng $(AB'D')$ có dạng $x+y+250z=0$.
		Vì $C'(1800;2500;-a)\in(AB'D')$, ta có: 
		 \[
		1800+2500-250a=0\Rightarrow  a=17{,}2.
		 \]
		Vậy  $a=17{,}2\,\text{ cm}$.
	} 
\end{bt}

% nguồn: Phan2_Lop_12_Hinh_hoc_Chuong_5
\begin{bt}%[2H5V1-3]%[ma:cc3d4cfb]
	Trong hệ trục tọa độ không gian $Oxyz$, cho $A(1;0;0)$, $B(0;b;0)$, $C(0;0;c)$, biết $b,c > 0$, phương trình mặt phẳng $(P)\colon y - z + 1 = 0$.  Tính $M = c + b$ biết $(ABC) \perp (P)$, $\mathrm{d}\left(O;(ABC)\right) = \dfrac{1}{3}$.
	\shortans{1} 
	 \loigiai{
		Phương trình mặt phẳng $(ABC)\colon \dfrac{x}{1} + \dfrac{y}{b}  + \dfrac{z}{c} = 1$.\\
		 Ta có $(ABC) \perp (P) \Rightarrow \dfrac{1}{b} - \dfrac{1}{c} = 0 \Leftrightarrow b=c$.
		$$\mathrm{d}(O;(ABC)) = \dfrac{|-1|}{\sqrt{1+\left(  \dfrac{1}{b}\right)^2+\left( \dfrac{1}{c}\right)^2}} = \dfrac{1}{3} \Leftrightarrow 9=1+2\left( \dfrac{1}{b}\right)^2, (b=c).$$
		Suy ra $b = \pm \dfrac{1}{2}$, do $b,c > 0$ nên $b=c= \dfrac{1}{2}$.  \\
		Vậy  $M=c+b=1$.
	} 
\end{bt}

% nguồn: Phan2_Lop_12_Hinh_hoc_Chuong_5
\begin{bt}%[2H5V1-3]%[ma:b2fc9c9e]
	Trong không gian $Oxyz$, cho điểm $B(6;4;0)$, $C(4;5;1)$, $M(2;1;6)$. Mặt phẳng  $(P)$ vuông góc với $BC$ và cách $M$ một khoảng bằng $\sqrt{6}$ có dạng $ax+y+cz+d=0$. Tính $-3a+c$.
	\shortans{7} 
	\loigiai{ 
		Mặt phẳng  $(P)$ vuông góc với $BC$ nên $\overrightarrow{n}_p = \overline{BC} = (-2;1;1) \Rightarrow (P)\colon -2x+y+z+D=0$.\\
		Lại  có $$\mathrm{d}(M;(P)) = \sqrt{6} \Leftrightarrow \dfrac{|D+3|}{\sqrt{6}} = \sqrt{6} \Leftrightarrow |D+3|=6 \Leftrightarrow
		\hoac{&D=-3 \\& D=3}$$ 
		Với  $D=-3 \Rightarrow (P)\colon -2x+y+z-3=0$.\\
		Với $D=3 \Rightarrow  (P)\colon -2x+y+z+3=0$.\\
		Vậy  $a=-2; c=1 \Rightarrow -3a+c=7$.
	 }
\end{bt}

% nguồn: Phan2_Lop_12_Hinh_hoc_Chuong_5
\begin{bt}%[2H5V1-3]%[ma:9e467937]
	Trong không gian $Oxyz$, mặt phẳng $(P)$ qua điểm $A(1;-2;-2)$, vuông góc với mặt phẳng $(Oxz)$ đồng thời khoảng cách từ điểm $B(3;1;-3)$ đến $(P)$ bằng $ \dfrac{3\sqrt{5}}{5}$. Phương trình mặt phẳng $(P)$  có dạng $2x + by + cz + d = 0$ ($d \neq 0$). Tính $P = 3b + 2c - d$.
	\shortans{46} 
	 \loigiai{
		Giả sử mặt phẳng cần lập có một vectơ pháp  tuyến là $\overrightarrow{n}_{(P)}=(a;b;c)$, $a^2+b^2+c^2 \neq 0$.\\
		Mặt phẳng $(P)$ vuông góc với mặt phẳng $(Oxz)\colon y=0$ nên $\overrightarrow{n}_{(P)} \cdot \overrightarrow{n}_{(Oxz)}  = 0 \Leftrightarrow b=0$.\\
		 $(P)$ đi qua điểm $A(1;-2;-2) \Rightarrow (P)\colon a(x-1) + c(z+2) = 0 \Leftrightarrow ax + cz - a + 2c = 0$.\\
		$d(B;(P)) = \dfrac{|2a-c|}{\sqrt{a^2+c^2}} = \dfrac{3\sqrt{5}}{5} \Leftrightarrow 5(2a-c)^2 = 9(a^2+c^2)$\\$\Leftrightarrow 11a^2  - 20ac - 4c^2 = 0 \Leftrightarrow \hoac {&a=2c \\&
		11a=-2c.}$\\ 
		Với $a=2c$ chọn $c=1  \Rightarrow a=2 \Rightarrow (P)\colon 2x+z=0$ (loại).\\
		Với $11a=-2c$ chọn $a=2 \Rightarrow c=-11 \Rightarrow  (P)\colon 2x-11z-24=0$ (thỏa mãn).\\
		Vậy $P=3b-2c-d =  0 - 2(-11) - (-24) = 46$.
	} 
\end{bt}

% nguồn: Phan2_Lop_12_Hinh_hoc_Chuong_5
\begin{bt}%[2H5V1-3]%[ma:a841431f]
	Trong không gian  với hệ tọa độ $Oxyz$, cho mặt phẳng $(\alpha)$ đi qua điểm $M(1; 2; 3)$ và cắt các tia $Ox$, $Oy$, $Oz$ tại $A, B, C$ sao cho độ dài $OA, OB, OC$ lập thành cấp số nhân có công bội bằng $3$. Phương trình mặt phẳng $(\alpha)$ có dạng $Ax + By + z + D = 0$, với $A, B, D \in \mathbb{R}$. Tính $A + B + D$.
	\shortans{-6} 
	\loigiai{ 
		Gọi $A(a; 0; 0)$, $B(0; b; 0)$, $C(0; 0; c)$ với $a,  b, c > 0$.\\
		Ta  có cấp số nhân $OA, OB, OC$ với công bội $3$ \\
		 Suy ra
		$$\heva{ 
		OB = 3 \cdot OA \Rightarrow b = 3a \\ 
		OC = 3 \cdot OB \Rightarrow c  = 9a
		 }$$
		Suy ra $A(a; 0; 0)$, $B(0;  3a; 0)$, $C(0; 0; 9a)$\\
		Phương trình mặt phẳng $(\alpha)$ theo đoạn chắn 
		$\dfrac{x}{a} + \dfrac{y}{3a} + \dfrac{z}{9a} =  1$\\
		Thay  $M(1; 2; 3)$ vào
		$$\dfrac{1}{a} + \dfrac{2}{3a} + \dfrac{3}{9a} = 1 \Rightarrow \dfrac{6}{a} = 3 \Rightarrow a = 2$$ 
		$$\Rightarrow (\alpha)\colon  \dfrac{x}{2} + \dfrac{y}{6} + \dfrac{z}{18} = 1 \Rightarrow 9x + 3y + z - 18 = 0.$$
		Suy ra  $A = 9$, $B = 3$, $D = -18 \Rightarrow A + B + D = -6.$
	 }
\end{bt}

% nguồn: Phan2_Lop_12_Hinh_hoc_Chuong_5
\begin{bt}%[2H5C1-3]%[ma:dee884ea]
	Một nắp bể nước hình chữ nhật $ABCD$ nằm cạnh bờ tường có kích thước $9\text{ dm}\times12\text{ dm}$ được kéo dài ra từ mặt sàn, do tác dụng của trọng lực nên nắp bể không thể mở ra được nếu không có người giữ. Người ta dùng một sợi dây xích dài $15\text{ dm}$ và kéo  căng nối đỉnh $C$ của hình chữ nhật với điểm $M$ nằm phía trên bờ tường sao cho $AM=9\text{ dm}$ và $AM$ vuông góc với mặt sàn. Chọn hệ trục $Oxyz$ như hình vẽ, khi đó nắp bể mở ra và tạo với mặt sàn một góc $\alpha$ (đơn vị trên mỗi trục tọa độ tính bằng dm). Bỏ qua độ dày của nắp bể.
	\begin{center}
		\begin{tikzpicture}[scale=0.5,>=stealth,font=\footnotesize,line join=round,line cap=round]
				 \draw[fill=gray!90] (-6.4,1.5)--(-6,1.3)--(-5.82,-3.49)--(-2.02,-6.02)--(5.8,-1.33)--(6.2,-1.6)--(-2,-6.52)--(-6.2,-3.72)--cycle;
				\draw[fill=gray!40] (-6.4,1.5)--(1.1,5.7)--(1.4,5.44)--(-6,1.3)--cycle; 
				 \draw[fill=gray!60] (-6,1.3)--(1.4,5.44)--(2,1.2)--(-5.82,-3.49)--cycle;
				\draw (1.4,5.44)--(2,1.2)--(6.2,-1.6); 
				\coordinate (A)  at (0,0);
				 \coordinate (C) at (-1.42,-2.71);
				\coordinate  (D) at (-4.3,-2.58);
				 \coordinate (B) at ($(A)+(C)-(D)$);
				\draw[fill=black!60,opacity=0.9] (A)--(B)--(C)--(D)--cycle; 
				\draw[fill=black!80]  (B)--(3.05,-0.28)--(-1.42,-2.97)--(-3.99,-2.85)--(-4.3,-2.58)--(-1.42,-2.71)--cycle;
				\coordinate (E) at (-1.73,-4.82); 
				\coordinate (F) at ($(A)+(E)-(D)$); 
				 \coordinate (M) at (0,3.72);
				\draw (D)--(E) 
					(A)--(F) 
					(E)--(F) 
					 (M) node[above right]{$M$}--(C);
				\fill[gray!40] (-5.82,-3.49)--(-2.02,-6.02)--(5.8,-1.33)--(2,1.2)--(A)--(B)--(3.05,-0.28)--(1.44,-1.25)--(F)--(E)--(D)--cycle; 
				\draw[->]  (0,0) node[above left,xshift=0.1cm,yshift=-0.1cm]{$O$}--(0,6) node[right]{$z$};
				\draw[->] (2,1.2)--(-7,-4.2) node[above]{$y$}; 
				\draw[->] (0,0)--(5.9,-3.5)  node[above]{$x$};
				\draw (A)  node[above right,yshift=0.2cm,xshift=-0.1cm]{$A$}
					(B) node[above]{$B$} 
					($(A)!0.5!(B)$)  node[above,xshift=0.1cm]{$9$ dm}
					($(A)!0.5!(D)$) node[above,xshift=-0.25cm]{$12$ dm} 
					(D)  node[above]{$D$}
					 (C) node[above left]{$C$}
					(D) node[xshift=0.7cm,yshift=-0.4cm]{$\alpha$}; 
			\draw[dashed] (C) to[bend left] (E);
		\draw[dashed] (B) to[bend left] (F);
		\draw[decorate,decoration={markings,mark=between positions 2mm and \pgfdecoratedpathlength-2mm step 2mm with{\draw (-3.5mm,-1.25mm)  to[out=0,in=160] (-2mm,-1.25mm) to[out=-20,in=160] (2mm,1.25mm) to[out=-20,in=180] (3.5mm,1.25mm);}}] (M) -- (C);
		\fill (A) circle (3pt); 
		\fill (B) circle (3pt); 
		 \fill (C) circle (3pt);
		 \fill (D) circle (3pt);
		\fill (M) circle (3pt); 
		\end{tikzpicture}
	\end{center}
	Phương trình mặt phẳng chứa nắp bể  sau khi kéo lên là $x+by+cz+d=0$. Tính $b+c^2-2d$.
	 \shortans{3}
	\loigiai{ 
		Điểm $M$ nằm trên mặt phẳng $(Oyz)$ với  $AM=9$ dm $\Rightarrow M(0;0;9)$.\\
		Chiếu $C$  lần lượt lên các mặt phẳng ta được $C(9\cos\alpha;\,12;\, 9\sin\alpha)$.\\
		Ta có $CM=15 \Leftrightarrow CM^2=15^2 \Leftrightarrow  (9\cos\alpha)^2+12^2+(9\sin\alpha-9)^2=15^2$\\
		$\Rightarrow  81\cos^2\alpha+144+81\sin^2\alpha-162\sin\alpha+81=225$\\
		 $\Rightarrow 81(\cos^2\alpha+\sin^2\alpha)-162\sin\alpha+225=225$\\
		$\Rightarrow -162\sin\alpha+81=0  \Rightarrow \sin\alpha=\dfrac{1}{2}\Rightarrow \alpha=30^\circ.$\\
		 Mặt phẳng kéo lên $(ABCD)$ qua điểm $A(0;0;0)$, chứa $AC$ và $AD$.\\
		Ta  có $\overrightarrow{AD}=(0;12;0)$ và $\overrightarrow{AC}=\left(\dfrac{9\sqrt3}{2};12;\dfrac{9}{2}\right)$.\\
		Véc-tơ  pháp tuyến của mặt $(ABCD)$ là
		$$\left[\overrightarrow{AC};\overrightarrow{AD}\right]=(54;0;-54\sqrt3)=54(1;0;-\sqrt3).$$ 
		Vậy phương trình mặt phẳng  chứa nắp bể sau khi kéo lên là
		$$x-\sqrt3z=0.$$ 
		 Suy ra $b=0$, $c=-\sqrt3$, $d=0$ $\Rightarrow b+c^2-2d=3.$\\
		Vậy  $b+c^2-2d=3.$
	 }
\end{bt}


%%%%%%%%%%%%%%%% TỰ LUẬN %%%%%%%%%%%%%%%%

% nguồn: Phan2_Lop_12_Hinh_hoc_Chuong_5
\begin{bt}%[2H5N1-3]%[ma:78c3b7c0]
	Trong không gian $Oxyz$. Viết  phương trình tổng quát của mặt phẳng $(P)$ đi qua điểm $M(2; 1;-3)$ có vectơ pháp tuyến là $\vec{n}=(3;-2; 6)$.
	 \loigiai{
		Phương trình tổng quát của mặt phẳng $(P)$ đi qua điểm $M(2; 1;-3)$ có vectơ pháp tuyến là $\vec{n}=(3;-2; 6)$ có  dạng: $3(x-2)-2(y-1)+6(z+3)=0\Leftrightarrow 3x-2y+6z+14=0$.
	} 
\end{bt}

% nguồn: Phan2_Lop_12_Hinh_hoc_Chuong_5
\begin{bt}%[2H5N1-3]%[ma:0a3940b8]
	Trong không gian $Oxyz$. Viết phương trình tổng quát của mặt phẳng $(P)$ đi qua điểm $M(2;-1; 0)$ và có cặp vectơ chỉ phương là  $\vec{a}=(2; 1; 3), \vec{b}=(1; 1; 2)$.
	 \loigiai{
		Mặt phẳng $(P)$ có cặp vectơ chỉ phương là $\vec{a}=(2;  1; 3), \vec{b}=(1; 1; 2)$ nên có vectơ pháp tuyến là $\vec{n}=\left[\vec{a}, \vec{b}\right]=(-1;-1; 1)$.\\
		Phương trình tổng quát của mặt phẳng $(P)$ đi qua điểm $M(2;-1; 0)$ có vectơ pháp tuyến là $\vec{n}=(-1;-1; 1)$ có dạng: $-1(x-2)-1(y+1)+1(z-0)=0\Leftrightarrow-x-y+z+1=0$. 
	 }
\end{bt}

% nguồn: Phan2_Lop_12_Hinh_hoc_Chuong_5
\begin{bt}%[2H5N1-3]%[ma:aa84de29]
	Lập phương trình mặt  phẳng theo đoạn chắn của mặt phẳng $(P)$, biết rằng mặt phẳng $(P)$ đi qua ba điểm $A(5; 0; 0)$, $B(0; 3; 0)$, $C(0; 0; 6)$.
	 \loigiai{
		Mặt phẳng  đi qua ba điểm $A(5; 0; 0)$, $B(0; 3; 0)$, $C(0; 0; 6)$ có phương trình là $\dfrac{x}{5}+\dfrac{y}{3}+\dfrac{z}{6}=1$.
	} 
\end{bt}

% nguồn: Phan2_Lop_12_Hinh_hoc_Chuong_5
\begin{bt}%[2H5N1-3]%[ma:058abb92]
	Trong không gian $Oxyz$, viết phương trình của mặt phẳng đi qua điểm $M(1; 2;-1)$ và vuông góc với trục $Ox$. 
	\loigiai{ 
		Mặt phẳng đi qua điểm $M(1; 2;-1)$ và vuông góc với trục $Ox$ thì nhận vectơ $\vec{i}=(1; 0; 0)$ làm vectơ pháp tuyến. Phương trình mặt  phẳng cần tìm là \[
		1(x-1)+0(y-2)+0(z+1)=0\Leftrightarrow  x-1=0.
		 \]
	} 
\end{bt}

% nguồn: Phan2_Lop_12_Hinh_hoc_Chuong_5
\begin{bt}%[2H5H1-3]%[ma:0601f966]
	 Trong không gian $Oxyz$. Viết phương trình mặt phẳng $(P)$ chứa điểm $M(1; 3;-2)$ cắt các tia $Ox$, $Oy$, $Oz$ lần lượt tại $A$, $B$, $C$ sao cho $\dfrac{OA}{1}=\dfrac{OB}{2}=\dfrac{OC}{4}$.
	\loigiai{ 
		Phương trình mặt chắn cắt tia $Ox$ tại $A(a; 0; 0)$, cắt tia $Oy$ tại $B(0; b; 0)$, cắt tia $Oz$ tại $C(0; 0; c)$ có dạng là $(P)\colon \dfrac{x}{a}+\dfrac{y}{b}+\dfrac{z}{c}=1$ (với $a > 0, b >  0, c > 0$).\\
		Theo đề $\dfrac{OA}{1}=\dfrac{OB}{2}=\dfrac{OC}{4} \Leftrightarrow \dfrac{a}{1}=\dfrac{b}{2}=\dfrac{c}{4} \Rightarrow\heva{&a=\dfrac{b}{2}  \\ &c=2b.}$\\
		Vì $M(1; 3;-2)$ nằm  trên mặt phẳng $(P)$ nên ta có $\dfrac{1}{\dfrac{b}{2}}+\dfrac{3}{b}+\dfrac{-2}{2b}=1\Leftrightarrow \dfrac{4}{b}=1\Leftrightarrow b=4$.\\
		Khi  đó $a=2, c=8$.\\
		Vậy phương  trình mặt phẳng $(P)$ là $\dfrac{x}{2}+\dfrac{y}{4}+\dfrac{z}{8}=1\Leftrightarrow 4x+2y+z-8=0$.
	} 
\end{bt}

% nguồn: Phan2_Lop_12_Hinh_hoc_Chuong_5
\begin{bt}%[2H5H1-3]%[ma:cd0c3f1d]
	Trong không gian $Oxyz$, viết phương trình mặt phẳng đi qua $M(2; 3;-1)$, song song với trục $Ox$ và vuông góc với  mặt phẳng $(Q)\colon x+2y-3z+1=0$.
	 \loigiai{
		Gọi  $(P)$ là mặt phẳng cần tìm.\\
		Do $(P)  \parallel Ox$ nên mặt phẳng $(P)$ có vectơ chỉ phương $\overrightarrow{u}_1=\vec{i}=(1; 0; 0)$.\\
		Mặt khác $(P) \perp(Q) \Rightarrow$ mặt phẳng $(P)$  có vectơ chỉ phương $\overrightarrow{u}_2=\overrightarrow{n}_{(Q)}=(1; 2;-3)$.\\
		Suy ra $ \overrightarrow{n}_{(P)}=\left[\overrightarrow{u}_1, \overrightarrow{u}_2\right]=(0;  3; 2)$.\\
		Mặt phẳng $(P)$ đi qua điểm $M(2; 3;-1)$ có phương trình $(P)\colon 3y+2z-7=0$. 
	 }
\end{bt}

% nguồn: Phan2_Lop_12_Hinh_hoc_Chuong_5
\begin{bt}%[2H5H1-3]%[ma:d027d58c]
	Trong không  gian với hệ trục $Oxyz$, cho mặt phẳng $(P)\colon 2x-2y+z-5=0$. Viết phương trình mặt phẳng $(Q)$ song song với mặt phẳng $(P)$, cách $(P)$ một khoảng bằng $3$ và cắt trục $Ox$ tại điểm có hoành độ dương.
	 \loigiai{
		Do $(Q) \parallel(P)$ nên mặt phẳng $(Q)\colon  2x-2y+z+C=0$; $C\neq-5$, chọn $M(0; 0; 5) \in(P)$.\\
		Ta  có $\mathrm{d}\left((P),(Q)\right)=\mathrm{d}\left(M,(Q)\right)=\dfrac{|5+C|}{\sqrt{2^2+(-2)^2+1^2}}=3\Leftrightarrow\hoac{&C=4\\&C=-14.}$\\
		Với $C=4\Rightarrow(Q)\colon 2x-2y+z+4=0$ khi đó $(Q)$ cắt $Ox$ tại điểm $M_1(-2; 0; 0)$ có hoành độ âm nên trường hợp này $(Q)$ không thỏa đề  bài.\\
		Với $C=-14\Rightarrow(Q)\colon 2x-2y+z-14=0$ khi đó $(Q)$ cắt $Ox$ tại điểm $M_2(7; 0; 0)$ có hoành  độ dương do đó $(Q)\colon 2x-2y+z-14=0$ thỏa mãn đề bài.\\
		Vậy phương  trình mặt phẳng $(Q)\colon 2x-2y+z-14=0$.
	} 
\end{bt}

% nguồn: Phan2_Lop_12_Hinh_hoc_Chuong_5
\begin{bt}%[2H5H1-3]%[ma:51398b9b]
	Trong không gian $Oxyz$, cho hai  điểm $A(4; 0; 1)$ và $B(-2; 2; 3)$. Viết phương trình mặt phẳng trung trực của đoạn thẳng $AB$.
	\loigiai{ 
		Mặt phẳng trung trực của đoạn thẳng $AB$ có vectơ pháp tuyến là $\overrightarrow{AB}=(-6; 2; 2)$ và đi qua trung điểm $I(1;  1; 2)$ của đoạn thẳng $AB$.\\
		Do  đó, ta có $-6(x-1)+2(y-1)+2(z-2)=0\Leftrightarrow-6x+2y+2z=0\Leftrightarrow 3x-y-z=0$.
	 }
\end{bt}

% nguồn: Phan2_Lop_12_Hinh_hoc_Chuong_5
\begin{bt}%[2H5H1-3]%[ma:efb265bc]
	Lập phương trình mặt phẳng $(P)$ đi qua ba điểm $A(1; 1; 1)$; $B(0; 4; 0)$;  $C(2; 2; 0)$.
	 \loigiai{
		 Ta có $\overrightarrow{AB}=(-1; 3;-1)$; $\overrightarrow{AC}=(1; 1;-1)$.\\
		 Xét \[
		\vec{n} = \left[\overrightarrow{AB}, \overrightarrow{AC}\right] = 
		 \left(
		\left|\begin{array}{cc}
		3 & -1  \\
		1  & -1
	\end{array}\right|,
	\left|\begin{array}{cc}
	-1 & -1 \\ 
	 -1 & 1
	\end{array}\right|,
	\left|\begin{array}{cc}
	-1 &  3 \\
	1 &  1
	\end{array}\right|
	 \right)
	=  (-2;\;-2;\;-4).
	 \]
	Khi đó $\vec{n}$ là một vectơ pháp  tuyến của mặt phẳng $(P)$.\\
	Mặt phẳng $(P)$ đi qua  $B(0; 4; 0)$ nên có phương trình
	\[-2(x-0)-2(y-4)-4(z-0)=0\Leftrightarrow  x+y+2z-4=0.\]
	Vậy  $(P)\colon x+y+2z-4=0$.
 }
\end{bt}

% nguồn: Phan2_Lop_12_Hinh_hoc_Chuong_5
\begin{bt}%[2H5H1-3]%[ma:cb9df7a5]
	Trong không gian $Oxyz$, cho hai điểm $A(-3; 2; 1)$ và $B(5;-4; 1)$. Gọi $M$ là hình chiếu vuông góc của $A$ trên $(Oxy)$, và $N$ là điểm đối xứng với $B$ qua $(Oy z)$. Viết phương trình mặt phẳng trung  trực $(P)$ của đoạn thẳng $MN$.
	\loigiai{ 
		$M$ là hình chiếu của $A(-3; 2; 1)$ trên trục $(Oxy)$ nên ta có $M(-3; 2;  0)$.
		$N$ là đối xứng với $B(5;-4; 1)$ qua $(Oy z)$ nên ta  có $N(-5;-4; 1)$.\\
		 Gọi $I$ là trung điểm $MN$.\\
		Ta có  $I\left(-4;-1; \dfrac{1}{2}\right)$ và $\overrightarrow{MN}=(-2;-6; 1)$.\\
		Phương trình mặt  phẳng $2(x+4)+6(y+1)-1\left(z-\dfrac{1}{2}\right)=0\Leftrightarrow 4x+12y-2z+29=0$.
	 }
\end{bt}

% nguồn: Phan2_Lop_12_Hinh_hoc_Chuong_5
\begin{bt}%[2H5H1-3]%[ma:253bf31b]
	Trong không gian với hệ trục tọa độ $Oxyz$. Viết phương trình mặt phẳng $(P)$ đi qua hai điểm $A(0; 1; 0)$, $B(2;  3; 1)$ và vuông góc với mặt phẳng $(Q)\colon x+2y-z=0$.
	 \loigiai{
		Ta có $\overrightarrow{AB}=(2; 2; 1)$, vectơ pháp tuyến mặt  phẳng $(Q)$ là $\overrightarrow{n}_{(Q)}=(1; 2;-1)$.\\
		Theo đề bài ta có vectơ pháp tuyến mặt  phẳng $(P)$ là $\overrightarrow{n}_{(P)}=\left[\overrightarrow{n}_{(Q)}, \overrightarrow{AB}\right]=(4;-3;-2)$.\\
		Phương trình mặt phẳng $(P)$ có dạng $4x-3y-2z+C=0$.\\ 
		Mặt phẳng $(P)$ đi  qua $A(0; 1; 0)$ nên $-3+C=0\Leftrightarrow C=3$.\\
		Vậy phương trình mặt phẳng $(P)$ là $4x-3y-2z+3=0$. 
	 }
\end{bt}

% nguồn: Phan2_Lop_12_Hinh_hoc_Chuong_5
\begin{bt}%[2H5H1-3]%[ma:0b47bd47]
	Trong không gian $Oxyz$, cho mặt phẳng $(P)$ đi qua giao tuyến của hai mặt phẳng $\left(\alpha_1\right)\colon 3x-y+z-2=0$, $\left(\alpha_2\right)\colon x+4y-5=0$ và  song song với mặt phẳng $\left(\alpha_3\right)\colon 2x+21y-z+7=0$. Viết phương trình của mặt phẳng $(P)$.
	 \loigiai{
		Ta có: $(P)  \parallel\left(\alpha_3\right) \Rightarrow(P)\colon 2x+21y-z+m=0~(m \neq-7)$.\\
		Gọi $d$ là giao tuyến  của $\left(\alpha_1\right)$; $\left(\alpha_2\right)$.\\ $d\colon\heva{&\left(\alpha_1\right)\colon 3x-y+z-2=0\\
		&\left(\alpha_2\right)\colon x+4y-5=0.}$\\ 
		Cho $y=0$ suy  ra $\heva{&3x+z-2=0\\ &x-5=0} \Rightarrow M(5; 0;-13)$.\\
		 $M\in(P) \Rightarrow 2\cdot 5+21\cdot 0-(-13)+m=0\Rightarrow m=-23$.\\
		Vậy  $2x+21y-z-23=0$.
	 }
\end{bt}

% nguồn: Phan2_Lop_12_Hinh_hoc_Chuong_5
\begin{bt}%[2H5H1-3]%[ma:50c92f2e]
	Trong không gian $Oxyz$, cho hình vuông $ABCD$ biết $A(1; 2; 1)$, $B(3;  0; 0)$, $C(1;-1;-2)$, $D(-1; 1;-1)$. Giả sử $I(a; b; c)$ là tâm đường tròn ngoại tiếp hình vuông $ABCD$ và $G$ là trọng tâm $\triangle ABC$. Viết phương trình mặt phẳng trung trực của $GI$.
	\loigiai{ 
		Vì $ABCD$ là hình vuông nên tâm  đường tròn ngoại tiếp hình vuông $ABCD$ là trung điểm của $AC$ và $BD$.\\
		Khi đó $I\left(1; \dfrac{1}{2};-\dfrac{1}{2}\right)$ và do $G$ là trọng tâm tam giác  $ABC$ nên $G\left(\dfrac{5}{3}; \dfrac{1}{3}; -\dfrac{1}{3}\right)$.\\
		Gọi $M$  là trung điểm của đoạn thẳng $GI$ nên $M\left(\dfrac{4}{3}; \dfrac{5}{12};-\dfrac{5}{12}\right)$.\\
		 Phương trình mặt phẳng trung trực đoạn thẳng $GI$ là
		$$-4\left(x-\dfrac{4}{3}\right)+\left(y-\dfrac{5}{12}\right)-\left(z+\dfrac{5}{12}\right)=0 \Leftrightarrow 8 x-2  y+2 z-9=0.
		 $$
	} 
\end{bt}

% nguồn: Phan2_Lop_12_Hinh_hoc_Chuong_5
\begin{bt}%[2H5H1-3]%[ma:526bda88]
	Lập phương trình mặt phẳng  $(P)$ đi qua
	\begin{enumerate}
		\item Điểm $I(3;-4;  1)$ và vuông góc với trục $Ox$.
		\item Điểm $K(-2; 4;-1)$ và  song song với mặt phẳng ($Ozx$).
		\item Điểm  $K(-2; 4;-1)$ và song song với mặt phẳng $(Q)\colon 3x+7y+10z+1=0$.
	\end{enumerate}
	\loigiai{ 
		\begin{enumerate}
			\item Mặt phẳng $(P) \perp Ox$ nên có một vectơ pháp tuyến là $\vec{n}_{(P)}=(1; 0;  0)$.\\
			Phương trình mặt phẳng $(P)$ là  $x-3=0$.
			\item Mặt phẳng  $(P) \parallel(Ozx) \Rightarrow(P) \perp Oy \Rightarrow \vec{n}_{(P)}=(0; 1; 0)$.\\
			Phương trình  mặt phẳng $(P)$ là $y-4=0$.
			\item Mặt phẳng $(P) \parallel(Q)  \Rightarrow \vec{n}_{(P)}=\vec{n}_{(Q)}=(3; 7; 10)$. \\
			 Phương trình mặt phẳng $(P)$ là $3(x+2)+7(y-4)+10(z+1)=0\Leftrightarrow 3x+7y+10z-12=0$.
		\end{enumerate}
	} 
\end{bt}

% nguồn: Phan2_Lop_12_Hinh_hoc_Chuong_5
\begin{bt}%[2H5H1-3]%[ma:12025cb6]
	Trong không gian với hệ tọa độ $Oxyz$, cho hình hộp $ABCD.A'B'C'D'$ có $A(1;-1; 3)$, $B(0;  2; 4)$ $D(2;-1; 1), A'(0; 1; 2)$.
	\begin{enumerate}
		\item  Tính tọa độ các điểm $C$, $B'$, $D'$.
		\item Viết phương  trình mặt phẳng $\left(CB' D'\right)$.
	\end{enumerate}
	 \loigiai{
		\begin{center}
			\begin{tikzpicture}[declare function={goc=225; gocx=75; a=3; b=8; h=5;},scale=0.7]
					\path 
						(0,0) coordinate  (A)
						(b,0)  coordinate (B)
						(goc:a)  coordinate (D)
						 ($(D)+(B)-(A)$) coordinate (C)
						(gocx:h) coordinate (A') 
						($(gocx:h)+(B)$) coordinate  (B')
						($(gocx:h)+(D)$)  coordinate (D')
						($(gocx:h)+(C)$)  coordinate (C');
					 \draw[dashed] (D)--(A)--(B)
						(A)--(A'); 
					\draw  (D)--(C)--(B)--(B')--(A')--(D')--(C')--(C)
						(D)--(D') 
						 (B')--(C');
					\foreach \x/\g in  {A/135,B/0,C/0,D/180,A'/135,B'/-30,C'/-15,D'/180}{\draw (\x) circle (1pt) node[shift={(\g:7pt)},font=\footnotesize]{$\x$};}
				\end{tikzpicture}
		\end{center}
		\begin{enumerate}
			\item Tính tọa độ các điểm $C$, $B'$,  $D'$.\\
			Ta  có $\overrightarrow{AD}=\overrightarrow{BC}\Leftrightarrow\heva{&x_C=1\\ &y_C-2=0\\&z_C-4=-2} \Leftrightarrow\heva{&x_C=1\\
			&y_C=2\\&z_C=2}$ suy ra $C(1;  2; 2)$.\\
			$\overrightarrow{AA'}=\overrightarrow{BB'}  \Leftrightarrow\heva{&x_{B'}=-1\\&y_{B'}-2=2\\ &z_{B'}-4=-1} \Leftrightarrow\heva{&x_B=-1\\&y_B=4\\&z_B=3}$ suy ra $B(-1; 4; 3)$.\\
			$\overrightarrow{AD}=\overrightarrow{A'D'} \Leftrightarrow\heva{&x_{D'}=1\\&y_{D'}-1=0\\&z_{D'}-2=-2}  \Leftrightarrow\heva{&x_{D'}=1\\&y_{D}=1\\ &z_{D'}=0}$ suy ra $D'(1; 1; 0)$.
			 \item Viết phương trình mặt phẳng $\left(CD' D'\right)$.\\
			Ta có $\overrightarrow{CB'}=(-2; 2; 1), \overrightarrow{CD'}=(0;-1;-2)$ là cặp vectơ chỉ phương của mặt phẳng $\left(CD' D'\right)$ nên $\left(CB' D'\right)$ có vectơ pháp tuyến là $\vec{n}=\left[\overrightarrow{CB'}, \overrightarrow{CD'}\right]=(-3;-4; 2)$ và đi qua điểm $C(1; 2;  2)$ có phương trình là $3x+4y-2z-7=0$.
		\end{enumerate}
	 }
\end{bt}

% nguồn: Phan2_Lop_12_Hinh_hoc_Chuong_5
\begin{bt}%[2H5H1-3]%[ma:a7bea2bb]
	Trong không gian $Oxyz$, viết phương trình mặt phẳng $(P)$ đi qua điểm $M(1;-1; 5)$ và vuông góc với hai mặt phẳng $(Q)\colon  3x+2y-z=0$, $(R)\colon x+y-z=0$.
	\loigiai{ 
		Mặt phẳng $(Q)$ và mặt  phẳng $(R)$ lần lượt có vectơ pháp tuyến là $\vec{n}_{(Q)}=(3; 2;-1), \vec{n}_{(R)}=(1; 1;-1)$.\\
		Vì mặt phẳng $(P)$ vuông góc với cả hai mặt phẳng $(R)$, $(Q)$ nên giá của các vectơ $\vec{n}_{(R)}$,  $\vec{n}_{(Q)}$ song song hoặc nằm trong mặt phẳng $(P)$.
		Suy ra $\left[\vec{n}_{(Q)}, \vec{n}_{(R)}\right]=(-1; 2; 1)$ là vectơ pháp tuyến của mặt phẳng  $(P)$.\\
		 Mặt phẳng $(P)$ đi qua điểm $M(1;-1; 5)$ và có vectơ pháp tuyến $\overrightarrow{n}_{(P)}=(-1; 2; 1)$ có phương trình là \[-1(x-1)+2(y+1)+1(z-5)=0\Leftrightarrow x-2y-z+2=0.\]
	} 
\end{bt}

% nguồn: Phan2_Lop_12_Hinh_hoc_Chuong_5
\begin{bt}%[2H5H1-3]%[ma:d7818505]
	Trong không gian $Oxyz$, cho hai mặt phẳng $(\alpha)\colon 3x-2y+2z+7=0$, $(\beta)\colon 5x-4y+3z+1=0$. Viết phương trình mặt phẳng đi qua gốc tọa độ $O$ đồng thời vuông góc với cả  $(\alpha)$ và $(\beta)$.
	 \loigiai{
		Vectơ pháp tuyến  của hai mặt phẳng lần lượt là $\vec{n}_{(\alpha)}=(3;-2; 2)$, $\vec{n}_{(\beta)}=(5;-4; 3)$.\\
		Khi đó vecto pháp  tuyến mặt phẳng cần tìm là $\vec{n}=\left[\vec{n}_{(\alpha)}, \vec{n}_{(\beta)}\right]=(2; 1;-2)$.\\
		Mặt phẳng đi qua gốc tọa độ  $O$ có vectơ pháp tuyến $\vec{n}=(2; 1;-2)$ là $2x+y-2z=0$.
	 }
\end{bt}

% nguồn: Ok_TuanLuan_Toan11_Lan2
\begin{ex}%[2H5V1-3]%[ma:0cbd491b]
	Trong không gian $Oxyz$, cho mặt phẳng $(P)\colon 2x-2y+z-5=0$. Viết phương trình mặt phẳng $(Q)$  song song với mặt phẳng $(P)$, cách $(P)$ một khoảng bằng $3$ và cắt trục $Ox$ tại điểm có hoành độ dương.
	 \loigiai{
		Ta có $(Q)\parallel (P)$ nên  mặt phẳng $(Q)\colon 2x-2y+z+C=0$; $C\ne-5$, chọn $M(0;0;5)\in(P)$.\\
		Ta có $\mathrm{d}\Big((P);(Q)\Big)=\mathrm{d}\Big(M;(Q)\Big)=\dfrac{\left|5+C\right|}{\sqrt{2^2+(-2)^2+1^2}}=3\Leftrightarrow \hoac{&C=4\\&C=-14.}$ 
		\begin{itemize}
			\item Với $C=4\Rightarrow (Q)\colon 2x-2y+z+4=0$ khi đó $(Q)$ cắt $Ox$ tại điểm $M_1(-2;0;0)$ có hoành độ âm  nên trường hợp này $(Q)$ không thỏa đề bài.
			\item Với $C=-14\Rightarrow (Q)\colon 2x-2y+z-14=0$ khi đó $(Q)$ cắt $Ox$ tại điểm $M_2(7;0;0)$  có hoành độ dương do đó $(Q)\colon 2x-2y+z-14=0$ thỏa đề bài.
		\end{itemize}
		Vậy phương trình mặt phẳng  $(Q)\colon 2x-2y+z-14=0$.
	} 
\end{ex}

% nguồn: Ok_TuanLuan_Toan11_Lan2
\begin{ex}%[2H5V1-3]%[ma:2394625f]
	Trong không gian $Oxyz$, cho mặt phẳng $(Q)\colon x+2y+2z-3=0$ và mặt phẳng $(P)$ không  qua $O$, song song mặt phẳng $(Q)$ và $\mathrm{d}\Big((P);(Q)\Big)=1$. Viết phương trình mặt phẳng $(P)$.
	\loigiai{ 
		Gọi phương trình mặt phẳng  $(P)$ có dạng $x+2y+2z+d=0$, với $d\ne 0$; $d\ne-3$.\\
		Có $\mathrm{d}\Big((P);(Q)\Big)=1\Leftrightarrow\dfrac{\left|d+3\right|}{\sqrt{1^2+2^2+2^2}}=1\Rightarrow  \hoac{&d=0\\&d=-6.}$\\
		Kết hợp điều kiện  ta suy ra $(P)$ có dạng $x+2y+2z-6=0$.
	 }
\end{ex}

% nguồn: Phan2_Lop_12_Hinh_hoc_Chuong_5
\begin{bt}%[2H5V1-3]%[ma:247be9b7]
	 Trong không gian $Oxyz$, cho mặt phẳng $(Q)\colon x+2y+2z-3=0$ và mặt phẳng $(P)$ không qua $O$, song song mặt phẳng $(Q)$ và $\mathrm{d}\Big((P);(Q)\Big)=1$. Viết phương trình mặt phẳng $(P)$.
	 \loigiai{
		Gọi phương trình  mặt phẳng $(P)$ có dạng $x+2y+2z+d=0$, với $d\ne 0$; $d\ne-3$.\\
		Có  $\mathrm{d}\Big((P);(Q)\Big)=1\Leftrightarrow\dfrac{\left|d+3\right|}{\sqrt{1^2+2^2+2^2}}=1\Rightarrow \hoac{&d=0\\&d=-6.}$\\
		Kết hợp điều kiện ta suy  ra $(P)$ có dạng $x+2y+2z-6=0$.
	 }
\end{bt}
